\input{preamble}

% OK, on commence ici.
%
\begin{document}

\title{Amorçage}

\maketitle

\phantomsection
\label{section-phantom}

\tableofcontents




\section{Introduction}
\label{section-introduction}

\noindent
Dans ce chapitre, nous utilisons les résultats des sections précédentes pour
donner des critères assurant qu'un préfaisceau d'ensembles sur la catégorie des
schémas est un espace algébrique. Une partie de ces résultats provient des
travaux d'Artin, voir \cite{ArtinI}, \cite{ArtinII},
\cite{Artin-Theorem-Representability},
\cite{Artin-Construction-Techniques},
\cite{Artin-Algebraic-Spaces},
\cite{Artin-Algebraic-Approximation},
\cite{Artin-Implicit-Function},
et \cite{ArtinVersal}.
Notre méthode consistera toutefois à employer autant que possible des arguments
semblables à ceux de l'article de Keel et Mori, voir
\cite{K-M}.

\section{Conventions}
\label{section-conventions}

\noindent
Nous supposons une fois pour toutes que tous les schémas considérés appartiennent
à un grand site fppf $\Sch_{fppf}$. De même, tout anneau $A$ considéré
est tel que $\Spec(A)$ est (isomorphe à) un
objet de ce grand site.

\medskip\noindent
Soient $S$ un schéma et $X$ un espace algébrique sur $S$.
Dans ce chapitre et dans les suivants, nous noterons $X \times_S X$
le produit de $X$ avec lui-même (dans la catégorie des espaces
algébriques sur $S$), plutôt que $X \times X$.




\section{Morphismes représentables par des espaces algébriques}
\label{section-morphism-representable-by-spaces}

\noindent
Nous définissons ici la notion de préfaisceau relativement représentable
par des espaces algébriques au-dessus d'un autre, puis en établissons quelques propriétés.

\begin{definition}
\label{definition-morphism-representable-by-spaces}
Soit $S$ un schéma appartenant à $\Sch_{fppf}$.
Soient $F$, $G$ des préfaisceaux sur $\Sch_{fppf}/S$.
On dit qu'un morphisme $a : F \to G$ est
{\it représentable par des espaces algébriques}
si, pour tout $U \in \Ob((\Sch/S)_{fppf})$ et
tout $\xi : U \to G$, le produit fibré $U \times_{\xi, G} F$
est un espace algébrique.
\end{definition}

\noindent
Voici une vérification élémentaire.

\begin{lemma}
\label{lemma-morphism-spaces-is-representable-by-spaces}
Soit $S$ un schéma.
Soit $f : X \to Y$ un morphisme d'espaces algébriques sur $S$.
Alors $f$ est représentable par des espaces algébriques.
\end{lemma}

\begin{proof}
C'est formel. Cela repose sur le fait que
la catégorie des espaces algébriques sur $S$ admet des produits fibrés, voir
Espaces, lemme \ref{spaces-lemma-fibre-product-spaces}.
\end{proof}

\begin{lemma}
\label{lemma-base-change-transformation}
\begin{slogan}
Tout changement de base d'un morphisme de préfaisceaux représentable
par des espaces algébriques est représentable par des espaces algébriques.
\end{slogan}
Soit $S$ un schéma. Soit
$$
\xymatrix{
G' \times_G F \ar[r] \ar[d]^{a'} & F \ar[d]^a \\
G' \ar[r] & G
}
$$
soit un carré cartésien de préfaisceaux sur $(\Sch/S)_{fppf}$.
Si $a$ est représentable par des espaces algébriques, alors $a'$ l'est aussi.
\end{lemma}

\begin{proof}
Omis. Indication : c'est formel.
\end{proof}

\begin{lemma}
\label{lemma-representable-by-spaces-transformation-to-sheaf}
Soit $S$ un schéma appartenant à $\Sch_{fppf}$.
Soient $F, G : (\Sch/S)_{fppf}^{opp} \to \textit{Sets}$.
Soit $a : F \to G$ représentable par des espaces algébriques.
Si $G$ est un faisceau, alors $F$ en est un aussi.
\end{lemma}

\begin{proof}
(Même démonstration que celle de
Espaces, lemme \ref{spaces-lemma-representable-transformation-to-sheaf}.)
Soit $\{\varphi_i : T_i \to T\}$ un recouvrement du site
$(\Sch/S)_{fppf}$.
Soient $s_i \in F(T_i)$ satisfaisant à la condition de faisceau.
Alors les $\sigma_i = a(s_i) \in G(T_i)$ satisfont eux aussi à la condition
de faisceau. Il existe donc un unique $\sigma \in G(T)$ tel
que $\sigma_i = \sigma|_{T_i}$. Par hypothèse,
$F' = h_T \times_{\sigma, G, a} F$ est un faisceau.
Notons que les $(\varphi_i, s_i) \in F'(T_i)$ satisfont eux aussi à la
condition de faisceau et proviennent donc d'un unique
$(\text{id}_T, s) \in F'(T)$. Manifestement, $s$ est la section de
$F$ cherchée.
\end{proof}

\begin{lemma}
\label{lemma-representable-by-spaces-transformation-diagonal}
Soit $S$ un schéma appartenant à $\Sch_{fppf}$.
Soient $F, G : (\Sch/S)_{fppf}^{opp} \to \textit{Sets}$.
Soit $a : F \to G$ représentable par des espaces algébriques.
Alors $\Delta_{F/G} : F \to F \times_G F$ est représentable par
des espaces algébriques.
\end{lemma}

\begin{proof}
(Même démonstration que celle de
Espaces, lemme \ref{spaces-lemma-representable-transformation-diagonal}.)
Soit $U$ un schéma. Soit $\xi = (\xi_1, \xi_2) \in (F \times_G F)(U)$.
Posons $\xi' = a(\xi_1) = a(\xi_2) \in G(U)$.
Par hypothèse, il existe un espace algébrique $V$ et un morphisme $V \to U$
qui représentent le produit fibré $U \times_{\xi', G} F$.
En particulier, les éléments $\xi_1, \xi_2$ donnent des morphismes
$f_1, f_2 : U \to V$ au-dessus de $U$. Comme $V$ représente le
produit fibré $U \times_{\xi', G} F$ et que
$\xi' = a \circ \xi_1 = a \circ \xi_2$,
on voit que, si $g : U' \to U$ est un morphisme, alors
$$
g^*\xi_1 = g^*\xi_2
\Leftrightarrow
f_1 \circ g = f_2 \circ g.
$$
Autrement dit, $U \times_{\xi, F \times_G F} F$
est représenté par $V \times_{\Delta, V \times V, (f_1, f_2)} U$,
qui est un espace algébrique.
\end{proof}

\noindent
La démonstration du
lemme \ref{lemma-representable-by-spaces-over-space}
ci-dessous est en fait quelque peu délicate. En effet,
nous ne pouvons employer l'argument de la démonstration de
Espaces, lemme \ref{spaces-lemma-representable-over-space},
car nous ne savons pas encore qu'un composé de transformations
représentables par des espaces algébriques est représentable par des espaces
algébriques. Nous utiliserons précisément ce lemme pour établir cet énoncé.

\begin{lemma}
\label{lemma-representable-by-spaces-over-space}
Soit $S$ un schéma appartenant à $\Sch_{fppf}$.
Soient $F, G : (\Sch/S)_{fppf}^{opp} \to \textit{Sets}$.
Soit $a : F \to G$ représentable par des espaces algébriques.
Si $G$ est un espace algébrique, alors $F$ en est un aussi.
\end{lemma}

\begin{proof}
Nous avons vu au
lemme \ref{lemma-representable-by-spaces-transformation-to-sheaf}
que $F$ est un faisceau.

\medskip\noindent
Soit $U$ un schéma et soit $U \to G$ un morphisme étale surjectif.
Alors $U \times_G F$ est un espace algébrique. Soit $W$ un schéma
et soit $W \to U \times_G F$ un morphisme étale surjectif.

\medskip\noindent
Affirmons d'abord que $W \to F$ est représentable.
Pour le voir, soient $X$ un schéma et $X \to F$ un morphisme.
Alors
$$
W \times_F X = W \times_{U \times_G F} U \times_G F \times_F X
= W \times_{U \times_G F} (U \times_G X)
$$
Comme $U \times_G F$ et $G$ sont tous deux des espaces algébriques, on voit
qu'il s'agit d'un schéma.

\medskip\noindent
Affirmons ensuite que $W \to F$ est étale et surjectif (cela a désormais
un sens puisque nous savons qu'il est représentable). Cela résulte de la
formule précédente : $W \to U \times_G F$ et $U \to G$
sont tous deux étales et surjectifs, donc
$W \times_{U \times_G F} (U \times_G X) \to U \times_G X$ and
$U \times_G X \to X$ sont étales et surjectifs, et le composé de
morphismes étales surjectifs est étale et surjectif.

\medskip\noindent
Posons $R = W \times_F W$. D'après ce qui précède, $R$ est un schéma et
les projections $t, s : R \to W$
sont étales. Il est clair que $R$ est une relation d'équivalence et que
$W \to F$ est un épimorphisme de faisceaux. Ainsi, $R$ est une relation
d'équivalence étale et $F = W/R$. Par conséquent, $F$ est un espace algébrique d'après
Espaces,
théorème \ref{spaces-theorem-presentation}.
\end{proof}

\begin{lemma}
\label{lemma-representable-by-spaces}
Soit $S$ un schéma.
Soit $a : F \to G$ un morphisme de préfaisceaux sur $(\Sch/S)_{fppf}$.
Supposons $a : F \to G$ représentable par des espaces algébriques.
Si $X$ est un espace algébrique sur $S$ et si $X \to G$ est un morphisme de préfaisceaux,
alors $X \times_G F$ est un espace algébrique.
\end{lemma}

\begin{proof}
D'après le lemme \ref{lemma-base-change-transformation}, la transformation
$X \times_G F \to X$ est représentable par des espaces algébriques. Sa source est donc
un espace algébrique d'après le
lemme \ref{lemma-representable-by-spaces-over-space}.
\end{proof}

\begin{lemma}
\label{lemma-composition-transformation}
Soit $S$ un schéma.
Soient
$$
\xymatrix{
F \ar[r]^a & G \ar[r]^b & H
}
$$
des morphismes de préfaisceaux sur $(\Sch/S)_{fppf}$.
Si $a$ et $b$ sont représentables par des espaces algébriques, alors
$b \circ a$ l'est aussi.
\end{lemma}

\begin{proof}
Soit $T$ un schéma sur $S$, et soit $T \to H$ un morphisme.
Par hypothèse, $T \times_H G$ est un espace algébrique. Le
lemme \ref{lemma-representable-by-spaces}
montre donc que $T \times_H F = (T \times_H G) \times_G F$ est lui aussi
un espace algébrique.
\end{proof}

\begin{lemma}
\label{lemma-product-transformations}
Soit $S$ un schéma.
Soient $F_i, G_i : (\Sch/S)_{fppf}^{opp} \to \textit{Sets}$, $i = 1, 2$.
Soient $a_i : F_i \to G_i$, $i = 1, 2$,
représentables par des espaces algébriques.
Alors
$$
a_1 \times a_2 : F_1 \times F_2 \longrightarrow G_1 \times G_2
$$
est représentable par des espaces algébriques.
\end{lemma}

\begin{proof}
Écrivons $a_1 \times a_2$ comme le composé
$F_1 \times F_2 \to G_1 \times F_2 \to G_1 \times G_2$.
La première flèche est le changement de base de $a_1$ par le morphisme
$G_1 \times F_2 \to G_1$, et la seconde
est le changement de base de $a_2$ par le morphisme
$G_1 \times G_2 \to G_2$. Ce lemme est donc une conséquence formelle des
lemmes \ref{lemma-composition-transformation}
et \ref{lemma-base-change-transformation}.
\end{proof}

\begin{lemma}
\label{lemma-representable-by-spaces-permanence}
Soit $S$ un schéma. Soient $a : F \to G$ et $b : G \to H$ des
transformations de foncteurs $(\Sch/S)_{fppf}^{opp} \to \textit{Sets}$.
Supposons que
\begin{enumerate}
\item $\Delta : G \to G \times_H G$ soit représentable
par des espaces algébriques, et que
\item $b \circ a : F \to H$ soit représentable par des espaces algébriques.
\end{enumerate}
Alors $a$ est représentable par des espaces algébriques.
\end{lemma}

\begin{proof}
Soit $U$ un schéma sur $S$ et soit $\xi \in G(U)$. Alors
$$
U \times_{\xi, G, a} F =
(U \times_{b(\xi), H, b \circ a} F) \times_{(\xi, a), (G \times_H G), \Delta} G
$$
Le résultat découle donc du lemme \ref{lemma-representable-by-spaces}.
\end{proof}

\begin{lemma}
\label{lemma-glueing-sheaves}
Soit $S \in \Ob(\Sch_{fppf})$. Soit $F$ un préfaisceau d'ensembles sur
$(\Sch/S)_{fppf}$. Supposons que
\begin{enumerate}
\item $F$ soit un faisceau pour la topologie de Zariski sur $(\Sch/S)_{fppf}$,
\item il existe un ensemble d'indices $I$ et des sous-foncteurs $F_i \subset F$ tels que
\begin{enumerate}
\item chaque $F_i$ soit un faisceau fppf,
\item chaque $F_i \to F$ soit représentable par des espaces algébriques,
\item $\coprod F_i \to F$ devienne surjectif après faisceautisation fppf.
\end{enumerate}
\end{enumerate}
Alors $F$ est un faisceau fppf.
\end{lemma}

\begin{proof}
Soient $T \in \Ob((\Sch/S)_{fppf})$ et $s \in F(T)$. D'après (2)(c),
il existe un recouvrement fppf $\{T_j \to T\}$ tel que
$s|_{T_j}$ soit une section de $F_{\alpha(j)}$ pour un certain $\alpha(j) \in I$.
Soit $W_j \subset T$ l'image de $T_j \to T$ ;
c'est un sous-schéma ouvert d'après Morphismes, lemme \ref{morphisms-lemma-fppf-open}.
D'après (2)(b),
$F_{\alpha(j)} \times_{F, s|_{W_j}} W_j \to W_j$ est un monomorphisme
d'espaces algébriques par lequel $T_j$ se factorise. Puisque $\{T_j \to W_j\}$
est un recouvrement fppf, on en déduit que
$F_{\alpha(j)} \times_{F, s|_{W_j}} W_j = W_j$, in other words
$s|_{W_j} \in F_{\alpha(j)}(W_j)$. On en conclut donc que
$\coprod F_i \to F$ est surjectif pour la topologie de Zariski.

\medskip\noindent
Soit $\{T_j \to T\}$ un recouvrement fppf de $(\Sch/S)_{fppf}$.
Soient $s, s' \in F(T)$ tels que $s|_{T_j} = s'|_{T_j}$ pour tout $j$.
Montrons que $s, s'$ sont égaux. Comme $F$ est un faisceau de Zariski d'après (1),
nous pouvons raisonner localement pour la topologie de Zariski sur $T$. D'après l'alinéa précédent,
nous pouvons supposer qu'il existe $i$ tel que $s \in F_i(T)$. Alors
$s'|_{T_j}$ est une section de $F_i$. D'après (2)(b),
$F_{i} \times_{F, s'} T \to T$ est un monomorphisme d'espaces algébriques
par lequel tous les $T_j$ se factorisent. On en conclut que
$s' \in F_i(T)$. Comme $F_i$ est un faisceau pour la topologie fppf,
on obtient $s = s'$.

\medskip\noindent
Soit $\{T_j \to T\}$ un recouvrement fppf de $(\Sch/S)_{fppf}$ et soient
$s_j \in F(T_j)$ tels que
$s_j|_{T_j \times_T T_{j'}} = s_{j'}|_{T_j \times_T T_{j'}}$. D'après l'hypothèse
(2)(c), nous pouvons raffiner le recouvrement et supposer que $s_j \in F_{\alpha(j)}(T_j)$
pour un certain $\alpha(j) \in I$. Soit $W_j \subset T$ l'image de $T_j \to T$,
qui est un sous-schéma ouvert d'après Morphismes, lemme \ref{morphisms-lemma-fppf-open}.
Alors $\{T_j \to W_j\}$ est un recouvrement fppf. Comme $F_{\alpha(j)}$ est un
sous-préfaisceau de $F$, les deux restrictions de $s_j$ à
$T_j \times_{W_j} T_j$ coïncident comme éléments de
$F_{\alpha(j)}(T_j \times_{W_j} T_j)$. La condition de faisceau pour
$F_{\alpha(j)}$ implique donc l'existence d'un $s'_j \in F_{\alpha(j)}(W_j)$
dont la restriction à $T_j$ est $s_j$. Pour toute paire d'indices
$j$ et $j'$, les sections $s'_j|_{W_j \cap W_{j'}}$ et
$s'_{j'}|_{W_j \cap W_{j'}}$ de $F$ coïncident d'après l'alinéa
précédent. Le fait que $F$ soit un faisceau de Zariski achève
la démonstration.
\end{proof}





\section{Propriétés des morphismes de préfaisceaux représentables par des espaces algébriques}
\label{section-representable-by-spaces-properties}

\noindent
Voici la définition qui convient.

\begin{definition}
\label{definition-property-transformation}
Soit $S$ un schéma. Soit $a : F \to G$ un morphisme de préfaisceaux sur
$(\Sch/S)_{fppf}$ représentable par des espaces algébriques.
Soit $\mathcal{P}$ une propriété des morphismes d'espaces algébriques qui
\begin{enumerate}
\item est stable par tout changement de base, et
\item est locale sur la base pour la topologie fppf, voir
Descente sur les espaces,
définition \ref{spaces-descent-definition-property-morphisms-local}.
\end{enumerate}
Dans ce cas, on dit que $a$ possède {\it la propriété $\mathcal{P}$} si, pour tout
schéma $U$ et tout $\xi : U \to G$, le morphisme d'espaces algébriques ainsi obtenu
$U \times_G F \to U$ possède la propriété $\mathcal{P}$.
\end{definition}

\noindent
Il importe de noter que nous n'utiliserons cette définition que pour
des propriétés de morphismes stables par changement de base et
locales sur la base pour la topologie fppf. Ce n'est
pas que la définition n'ait autrement aucun sens ; c'est plutôt que
nous pouvons vouloir donner une autre définition, mieux
adaptée à la propriété envisagée.

\medskip\noindent
La définition précédente s'applique\footnote{La stabilité par changement
de base résulte de
Morphismes d'espaces, lemmes
\ref{spaces-morphisms-lemma-base-change-surjective},
\ref{spaces-morphisms-lemma-base-change-quasi-compact},
\ref{spaces-morphisms-lemma-base-change-etale},
\ref{spaces-morphisms-lemma-base-change-smooth},
\ref{spaces-morphisms-lemma-base-change-flat},
\ref{spaces-morphisms-lemma-base-change-separated},
\ref{spaces-morphisms-lemma-base-change-finite-type},
\ref{spaces-morphisms-lemma-base-change-quasi-finite},
\ref{spaces-morphisms-lemma-base-change-finite-presentation},
\ref{spaces-morphisms-lemma-base-change-affine},
\ref{spaces-morphisms-lemma-base-change-proper}, and
Espaces, lemme
\ref{spaces-lemma-base-change-immersions}.
Le caractère local sur la base pour la topologie fppf résulte de
Descente sur les espaces, lemmes
\ref{spaces-descent-lemma-descending-property-surjective},
\ref{spaces-descent-lemma-descending-property-quasi-compact},
\ref{spaces-descent-lemma-descending-property-etale},
\ref{spaces-descent-lemma-descending-property-smooth},
\ref{spaces-descent-lemma-descending-property-flat},
\ref{spaces-descent-lemma-descending-property-separated},
\ref{spaces-descent-lemma-descending-property-finite-type},
\ref{spaces-descent-lemma-descending-property-quasi-finite},
\ref{spaces-descent-lemma-descending-property-locally-finite-presentation},
\ref{spaces-descent-lemma-descending-property-affine},
\ref{spaces-descent-lemma-descending-property-proper}, and
\ref{spaces-descent-lemma-descending-property-closed-immersion}.
}
par exemple aux propriétés d'être
« surjectif »,
« quasi-compact »,
« étale »,
« lisse »,
« plat »,
« séparé »,
« (localement) de type fini »,
« (localement) quasi-fini »,
« (localement) de présentation finie »,
« affine »,
« propre » et
« une immersion fermée ».
Autrement dit, $a$ est
{\it surjectif}
(resp.\ {\it quasi-compact},
{\it étale},
{\it lisse},
{\it plat},
{\it séparé},
{\it (localement) de type fini},
{\it (localement) quasi-fini},
{\it (localement) de présentation finie},
{\it propre},
{\it une immersion fermée})
si, pour tout schéma $T$ et tout morphisme $\xi : T \to G$,
le morphisme d'espaces algébriques $T \times_{\xi, G} F \to T$ est
surjectif
(resp.\ quasi-compact,
étale,
plat,
séparé,
(localement) de type fini,
(localement) quasi-fini,
(localement) de présentation finie,
propre,
une immersion fermée).

\medskip\noindent
Vérifions ensuite la compatibilité avec les notions déjà définies. D'après le
lemme \ref{lemma-morphism-spaces-is-representable-by-spaces},
tout morphisme entre espaces algébriques sur $S$ est représentable par
des espaces algébriques. D'autre part, d'après
Morphismes d'espaces,
lemme \ref{spaces-morphisms-lemma-surjective-local}
(resp.\ \ref{spaces-morphisms-lemma-quasi-compact-local},
\ref{spaces-morphisms-lemma-etale-local},
\ref{spaces-morphisms-lemma-smooth-local},
\ref{spaces-morphisms-lemma-flat-local},
\ref{spaces-morphisms-lemma-separated-local},
\ref{spaces-morphisms-lemma-finite-type-local},
\ref{spaces-morphisms-lemma-quasi-finite-local},
\ref{spaces-morphisms-lemma-finite-presentation-local},
\ref{spaces-morphisms-lemma-affine-local},
\ref{spaces-morphisms-lemma-proper-local},
\ref{spaces-morphisms-lemma-closed-immersion-local})
la définition précédente de
surjectif
(resp.\ quasi-compact,
étale,
lisse,
plat,
séparé,
(localement) de type fini,
(localement) quasi-fini,
(localement) de présentation finie,
affine,
propre,
immersion fermée)
coïncide avec la définition déjà donnée pour les morphismes
d'espaces algébriques.

\medskip\noindent
Voici maintenant quelques lemmes formels.

\begin{lemma}
\label{lemma-base-change-transformation-property}
Soit $S$ un schéma.
Soit $\mathcal{P}$ une propriété comme dans la
définition \ref{definition-property-transformation}.
Soit
$$
\xymatrix{
G' \times_G F \ar[r] \ar[d]^{a'} & F \ar[d]^a \\
G' \ar[r] & G
}
$$
un carré cartésien de préfaisceaux sur $(\Sch/S)_{fppf}$.
Si $a$ est représentable par des espaces algébriques et possède $\mathcal{P}$,
alors $a'$ aussi.
\end{lemma}

\begin{proof}
Omis. Indication : c'est formel.
\end{proof}

\begin{lemma}
\label{lemma-composition-transformation-property}
Soit $S$ un schéma.
Soit $\mathcal{P}$ une propriété comme dans la
définition \ref{definition-property-transformation},
et supposons $\mathcal{P}$ stable par composition.
Soient
$$
\xymatrix{
F \ar[r]^a & G \ar[r]^b & H
}
$$
des morphismes de préfaisceaux sur $(\Sch/S)_{fppf}$.
Si $a$ et $b$ sont représentables par des espaces algébriques et possèdent
$\mathcal{P}$, alors $b \circ a$ la possède aussi.
\end{lemma}

\begin{proof}
Omis. Indication : voir le
lemme \ref{lemma-composition-transformation}
et utiliser la stabilité par composition.
\end{proof}

\begin{lemma}
\label{lemma-product-transformations-property}
Soit $S$ un schéma.
Soient $F_i, G_i : (\Sch/S)_{fppf}^{opp} \to \textit{Sets}$,
$i = 1, 2$.
Soient $a_i : F_i \to G_i$, $i = 1, 2$, représentables par des espaces algébriques.
Soit $\mathcal{P}$ une propriété comme dans la
définition \ref{definition-property-transformation},
stable par composition.
Si $a_1$ et $a_2$ possèdent la propriété $\mathcal{P}$, alors
$a_1 \times a_2 : F_1 \times F_2 \longrightarrow G_1 \times G_2$ la possède aussi.
\end{lemma}

\begin{proof}
Notons que l'énoncé a un sens d'après le
lemme \ref{lemma-product-transformations}.
Démonstration omise.
\end{proof}

\begin{lemma}
\label{lemma-transformations-property-implication}
Soit $S$ un schéma.
Soient $F, G : (\Sch/S)_{fppf}^{opp} \to \textit{Sets}$.
Soit $a : F \to G$ une transformation de foncteurs représentable par
des espaces algébriques.
Soient $\mathcal{P}$, $\mathcal{P}'$ des propriétés comme dans la
définition \ref{definition-property-transformation}.
Supposons que, pour tout morphisme $f : X \to Y$ d'espaces algébriques sur $S$,
on ait $\mathcal{P}(f) \Rightarrow \mathcal{P}'(f)$.
Si $a$ possède la propriété $\mathcal{P}$, alors
$a$ possède la propriété $\mathcal{P}'$.
\end{lemma}

\begin{proof}
C'est formel.
\end{proof}

\begin{lemma}
\label{lemma-surjective-flat-locally-finite-presentation}
Soit $S$ un schéma.
Soient $F, G : (\Sch/S)_{fppf}^{opp} \to \textit{Sets}$ des faisceaux.
Soit $a : F \to G$ représentable par des espaces algébriques, plat,
localement de présentation finie et surjectif.
Alors $a : F \to G$ est surjectif comme morphisme de faisceaux.
\end{lemma}

\begin{proof}
Soit $T$ un schéma sur $S$ et soit $g : T \to G$ un point à valeurs dans
$T$ de $G$. Par hypothèse, $T' = F \times_G T$ est un espace algébrique et
le morphisme $T' \to T$ est un morphisme d'espaces algébriques plat,
localement de présentation finie et surjectif.
Soit $U \to T'$ un morphisme étale surjectif, où $U$ est un schéma.
Par définition des morphismes plats d'espaces algébriques,
le morphisme de schémas $U \to T$ est alors plat. Il en va de même pour la
propriété « localement de présentation finie ». Le morphisme $U \to T$ est aussi surjectif,
voir
Morphismes d'espaces, lemme \ref{spaces-morphisms-lemma-surjective-local}.
Ainsi, $\{U \to T\}$ est un recouvrement fppf tel
que $g|_U \in G(U)$ provienne d'un élément de $F(U)$, à savoir du
morphisme $U \to T' \to F$. Cela prouve que le morphisme est surjectif comme
morphisme de faisceaux, voir
Sites, définition \ref{sites-definition-sheaves-injective-surjective}.
\end{proof}




\section{Amorçage de la diagonale}
\label{section-bootstrap-diagonal}

\noindent
Dans cette section, nous prouvons que la diagonale d'un faisceau $F$ sur
$(\Sch/S)_{fppf}$ est représentable dès qu'il existe
un « recouvrement fppf » de $F$ par un schéma ou par un espace algébrique, voir le
lemme \ref{lemma-bootstrap-diagonal}.

\begin{lemma}
\label{lemma-representable-diagonal}
\begin{slogan}
La diagonale d'un préfaisceau est représentable par des espaces algébriques si et seulement si
tout morphisme d'un schéma vers ce préfaisceau est représentable par des espaces algébriques.
\end{slogan}
Soit $S$ un schéma.
Soit $F$ un préfaisceau sur $(\Sch/S)_{fppf}$.
Les conditions suivantes sont équivalentes :
\begin{enumerate}
\item $\Delta_F : F \to F \times F$ est représentable par des espaces algébriques ;
\item pour tout schéma $T$, tout morphisme $T \to F$ est représentable par des espaces
algébriques ;
\item pour tout espace algébrique $X$, tout morphisme $X \to F$ est représentable
par des espaces algébriques.
\end{enumerate}
\end{lemma}

\begin{proof}
Supposons (1). Soit $X \to F$ comme en (3). Soit $T$ un schéma et soit
$T \to F$ un morphisme. On a alors
$$
T \times_F X = (T \times_S X) \times_{F \times F, \Delta} F
$$
qui est un espace algébrique d'après le
lemme \ref{lemma-representable-by-spaces}
et (1). Ainsi, $X \to F$ est représentable, c'est-à-dire que (3) est vérifiée.
L'implication (3) $\Rightarrow$ (2) est immédiate. Supposons (2).
Soit $T$ un schéma et soit $(a, b) : T \to F \times F$ un morphisme.
Alors
$$
F \times_{\Delta_F, F \times F} T =
(T \times_{a, F, b} T) \times_{T \times T, \Delta_T} T
$$
qui est un espace algébrique par hypothèse. Ainsi, $\Delta_F$ est
représentable par des espaces algébriques, c'est-à-dire que (1) est vérifiée.
\end{proof}

\noindent
En particulier, si $F$ est un préfaisceau satisfaisant aux conditions équivalentes du
lemme, alors, pour tout morphisme $X \to F$ où $X$ est un espace algébrique,
il est légitime de dire que $X \to F$ est surjectif (resp.\ étale, plat,
localement de présentation finie) en appliquant la
définition \ref{definition-property-transformation}.

\medskip\noindent
Avant de procéder à l'amorçage proprement dit, établissons un petit lemme utile.

\begin{lemma}
\label{lemma-after-fppf-sep-lqf}
Soit $S$ un schéma.
Soit
$$
\xymatrix{
E \ar[r]_a \ar[d]_f & F \ar[d]^g \\
H \ar[r]^b & G
}
$$
un diagramme cartésien de faisceaux sur $(\Sch/S)_{fppf}$, de sorte que
$E = H \times_G F$. Si
\begin{enumerate}
\item $g$ est représentable par des espaces algébriques, surjectif, plat et
localement de présentation finie ;
\item $a$ est représentable par des espaces algébriques, séparé et
localement quasi-fini,
\end{enumerate}
alors $b$ est représentable (par des schémas), séparé et
localement quasi-fini.
\end{lemma}

\begin{proof}
Soit $T$ un schéma et soit $T \to G$ un morphisme.
Il faut montrer que $T \times_G H$ est un schéma et que
le morphisme $T \times_G H \to T$ est séparé et
localement quasi-fini. Nous pouvons donc effectuer sur tout le diagramme le changement de base vers $T$
et supposer que $G$ est un schéma. Alors $F$ est un espace algébrique.
Soit $U$ un schéma et soit $U \to F$ un morphisme étale surjectif.
Alors $U \to F$ est représentable, surjectif, plat et
localement de présentation finie d'après
Morphismes d'espaces,
lemmes \ref{spaces-morphisms-lemma-etale-flat} et
\ref{spaces-morphisms-lemma-etale-locally-finite-presentation}.
Le
lemme \ref{lemma-composition-transformation}
montre que $U \to G$ est lui aussi surjectif, plat et localement de présentation finie.
Notons que le changement de base $E \times_F U \to U$ de $a$ reste
séparé et localement quasi-fini (d'après le
lemme \ref{lemma-base-change-transformation-property}). Nous pouvons donc
remplacer la partie supérieure du diagramme du lemme par
$E \times_F U \to U$. Autrement dit, nous pouvons supposer que
$F \to G$ est un morphisme de schémas surjectif et plat,
localement de présentation finie.
En particulier, $\{F \to G\}$ est un recouvrement fppf de schémas.
D'après
Morphismes d'espaces, proposition
\ref{spaces-morphisms-proposition-locally-quasi-finite-separated-over-scheme}
on en conclut que $E$ est aussi un schéma.
D'après
Descente, lemme \ref{descent-lemma-descent-data-sheaves},
l'égalité $E = H \times_G F$ signifie que l'on obtient une donnée de descente
sur $E$ relativement au recouvrement fppf $\{F \to G\}$.
D'après
Compléments sur les morphismes, lemme
\ref{more-morphisms-lemma-separated-locally-quasi-finite-morphisms-fppf-descend}
Cette donnée de descente est effective.
En appliquant de nouveau
Descente, lemme \ref{descent-lemma-descent-data-sheaves},
on en déduit que $H$ est un schéma.
Enfin, d'après
Descente, lemmes \ref{descent-lemma-descending-property-separated} et
\ref{descent-lemma-descending-property-quasi-finite}
il s'ensuit que $b$ est séparé et localement quasi-fini.
\end{proof}

\noindent
Voici le résultat annoncé par le titre de la section.

\begin{lemma}
\label{lemma-bootstrap-diagonal}
Soit $S$ un schéma.
Soit $F : (\Sch/S)_{fppf}^{opp} \to \textit{Sets}$ un foncteur.
Supposons que
\begin{enumerate}
\item le préfaisceau $F$ soit un faisceau ;
\item il existe un espace algébrique $X$ et un morphisme $X \to F$
représentable par des espaces algébriques, surjectif, plat et
localement de présentation finie.
\end{enumerate}
Alors $\Delta_F$ est représentable (par des schémas).
\end{lemma}

\begin{proof}
Soit $U \to X$ un morphisme étale surjectif d'un schéma vers $X$.
Alors $U \to X$ est représentable, surjectif, plat et
localement de présentation finie d'après
Morphismes d'espaces,
lemmes \ref{spaces-morphisms-lemma-etale-flat} et
\ref{spaces-morphisms-lemma-etale-locally-finite-presentation}.
Le
lemme \ref{lemma-composition-transformation-property}
montre que le composé $U \to F$ est lui aussi représentable par des espaces algébriques,
surjectif, plat et localement de présentation finie.
Ainsi, $R = U \times_F U$ est un espace algébrique, voir le
lemme \ref{lemma-representable-by-spaces}.
Le morphisme d'espaces algébriques $R \to U \times_S U$ est
un monomorphisme, donc il est séparé (puisque la diagonale d'un monomorphisme
est un isomorphisme, voir
Morphismes d'espaces,
lemme \ref{spaces-morphisms-lemma-monomorphism}).
Comme $U \to F$ est localement de présentation finie, les deux
morphismes $R \to U$ le sont aussi, voir le
lemme \ref{lemma-base-change-transformation-property}.
Ainsi, $R \to U \times_S U$ est localement de type fini (utiliser
Morphismes d'espaces,
lemmes \ref{spaces-morphisms-lemma-finite-presentation-finite-type} et
\ref{spaces-morphisms-lemma-permanence-finite-type}).
Au total,
$R \to U \times_S U$ est un monomorphisme localement de type
fini, donc un morphisme séparé et localement quasi-fini, voir
Morphismes d'espaces, lemme
\ref{spaces-morphisms-lemma-monomorphism-loc-finite-type-loc-quasi-finite}.

\medskip\noindent
Nous sommes maintenant en mesure de prouver que $\Delta_F$ est représentable.
Soit $T$ un schéma et soit $(a, b) : T \to F \times F$ un morphisme.
Posons
$$
T' = (U \times_S U) \times_{F \times F} T.
$$
Notons que $U \times_S U \to F \times F$ est
représentable par des espaces algébriques, surjectif, plat et
localement de présentation finie d'après le
lemme \ref{lemma-product-transformations-property}.
Ainsi, $T'$ est un espace algébrique et le morphisme de projection
$T' \to T$ est surjectif, plat et localement de présentation finie.
Considérons $Z = T \times_{F \times F} F$ (c'est un faisceau) et
$$
Z' = T' \times_{U \times_S U} R
= T' \times_T Z.
$$
On voit que $Z'$ est un espace algébrique et que
$Z' \to T'$ est séparé et localement quasi-fini d'après la
discussion du premier alinéa de la démonstration, qui a établi que $R$ est
un espace algébrique et que le
morphisme $R \to U \times_S U$ possède ces propriétés.
Nous pouvons donc appliquer le
lemme \ref{lemma-after-fppf-sep-lqf}
au diagramme
$$
\xymatrix{
Z' \ar[r] \ar[d] & T' \ar[d] \\
Z \ar[r] & T
}
$$
ce qui permet de conclure.
\end{proof}

\noindent
Voici une variante du résultat précédent.

\begin{lemma}
\label{lemma-bootstrap-locally-quasi-finite}
Soit $S$ un schéma. Soit $F : (\Sch/S)_{fppf}^{opp} \to \textit{Sets}$ un
foncteur. Soit $X$ un schéma et soit $X \to F$ représentable par des espaces
algébriques et localement quasi-fini. Alors $X \to F$ est représentable
(par des schémas).
\end{lemma}

\begin{proof}
Soit $T$ un schéma et soit $T \to F$ un morphisme. Il faut montrer que
l'espace algébrique $X \times_F T$ est représentable par un schéma. Considérons
le morphisme
$$
X \times_F T  \longrightarrow X \times_{\Spec(\mathbf{Z})} T
$$
Comme $X \times_F T \to T$ est localement quasi-fini, la flèche affichée
l'est aussi (Morphismes d'espaces, lemme
\ref{spaces-morphisms-lemma-permanence-quasi-finite}).
D'autre part, la flèche affichée est un monomorphisme,
donc elle est séparée (Morphismes d'espaces, lemme
\ref{spaces-morphisms-lemma-monomorphism-separated}).
Ainsi, $X \times_F T$ est un schéma d'après Morphismes d'espaces, proposition
\ref{spaces-morphisms-proposition-locally-quasi-finite-separated-over-scheme}.
\end{proof}











\section{Amorçage}
\label{section-bootstrap}

\noindent
Prévenons d'emblée le lecteur que le résultat de cette section sera
supplanté par le théorème plus fort
\ref{theorem-final-bootstrap}.
En revanche, le théorème de cette section est nettement plus facile à
démontrer et donne déjà une bonne compréhension du mécanisme en jeu,
notamment aux lecteurs qui s'intéressent surtout aux champs de
Deligne--Mumford.

\medskip\noindent
Dans
Espaces, section \ref{spaces-section-algebraic-spaces},
nous avons défini un espace algébrique comme un faisceau pour la topologie fppf
dont la diagonale est représentable et vers lequel il existe un morphisme étale
surjectif provenant d'un schéma. Dans cette section, nous montrons qu'un
faisceau pour la topologie fppf dont la diagonale est représentable par des espaces
algébriques et qui admet un recouvrement étale surjectif par un espace algébrique
est lui aussi un espace algébrique.
Autrement dit, la catégorie des espaces algébriques s'obtient en élargissant la
catégorie des schémas par les faisceaux fppf $F$ qui possèdent une diagonale
représentable et un recouvrement étale par un schéma. Le
résultat de cette section affirme que recommencer ce procédé à partir de
la catégorie des espaces algébriques ne conduit pas à une nouvelle catégorie.

\medskip\noindent
Une autre motivation des résultats de cette section est qu'ils permettront
plus loin d'assurer qu'un champ de Deligne--Mumford dont le champ d'inertie est
trivial est équivalent à un espace algébrique, voir
Champs algébriques, lemme \ref{algebraic-lemma-algebraic-stack-no-automorphisms}.

\medskip\noindent
Voici le résultat principal de cette section (comme indiqué plus haut, il
sera supplanté par le théorème plus fort
\ref{theorem-final-bootstrap}).

\begin{theorem}
\label{theorem-bootstrap}
Soit $S$ un schéma.
Soit $F : (\Sch/S)_{fppf}^{opp} \to \textit{Sets}$ un foncteur.
Supposons que
\begin{enumerate}
\item le préfaisceau $F$ soit un faisceau ;
\item le morphisme diagonal $F  \to F \times F$ soit représentable par
des espaces algébriques ;
\item il existe un espace algébrique $X$
et un morphisme $X \to F$ surjectif et étale.
\end{enumerate}
ou bien supposons que
\begin{enumerate}
\item[(a)] le préfaisceau $F$ soit un faisceau ;
\item[(b)] il existe un espace algébrique $X$ et un morphisme $X \to F$
représentable par des espaces algébriques, surjectif et étale.
\end{enumerate}
Alors $F$ est un espace algébrique.
\end{theorem}

\begin{proof}
Nous utiliserons sans plus les mentionner les remarques qui suivent immédiatement la
définition \ref{definition-property-transformation}.

\medskip\noindent
Supposons (1), (2) et (3), et soit $X \to F$ comme en (3).
D'après le lemme \ref{lemma-representable-diagonal}, le morphisme
$X \to F$ est représentable par des espaces algébriques. Ainsi,
(a) et (b) sont vérifiées.

\medskip\noindent
Supposons (a) et (b), et soit $X \to F$ comme en (b).
Soit $U \to X$ un morphisme étale surjectif d'un schéma vers $X$.
D'après le lemme \ref{lemma-composition-transformation}, la transformation
$U \to F$ est représentable par des espaces algébriques, surjective et étale.
Pour prouver que $F$ est un espace algébrique, il suffit donc de montrer que
$\Delta_F$ est représentable (Espaces, définition
\ref{spaces-definition-algebraic-space}). Cela résulte immédiatement du
lemme \ref{lemma-bootstrap-diagonal}.
Nous pouvons aussi éviter ce lemme et montrer directement que $F$
est un espace algébrique, comme dans l'alinéa suivant.

\medskip\noindent
En effet, soient $U$ un schéma et $U \to F$ un morphisme représentable par des espaces
algébriques, surjectif et étale. Considérons le produit fibré $R = U \times_F U$.
Les deux projections $R \to U$ sont représentables par des espaces algébriques, surjectives
et étales (lemme \ref{lemma-base-change-transformation-property}).
En particulier, $R$ est un espace algébrique d'après le
lemme \ref{lemma-representable-by-spaces-over-space}.
Le morphisme d'espaces algébriques $R \to U \times_S U$ est un monomorphisme,
donc il est séparé (puisque la diagonale d'un monomorphisme est un isomorphisme).
Comme $R \to U$ est étale, $R \to U$ est localement quasi-fini, voir
Morphismes d'espaces,
lemme \ref{spaces-morphisms-lemma-etale-locally-quasi-finite}.
On en déduit que $R \to U \times_S U$ est lui aussi
localement quasi-fini d'après
Morphismes d'espaces,
lemme \ref{spaces-morphisms-lemma-permanence-quasi-finite}.
Ainsi,
Morphismes d'espaces, proposition
\ref{spaces-morphisms-proposition-locally-quasi-finite-separated-over-scheme}
s'applique et $R$ est un schéma. D'après le
lemme \ref{lemma-surjective-flat-locally-finite-presentation},
le morphisme $U \to F$ est un épimorphisme de faisceaux. Ainsi $F = U/R$.
On en conclut que $F$ est un espace algébrique d'après
Espaces, théorème \ref{spaces-theorem-presentation}.
\end{proof}









\section{Recherche d'ouverts}
\label{section-finding-opens}


\medskip\noindent
Commençons par établir un lemme qui améliore et généralise légèrement
Espaces, lemme \ref{spaces-lemma-finding-opens},
aux faisceaux quotients associés à des groupoïdes.

\begin{lemma}
\label{lemma-better-finding-opens}
Soit $S$ un schéma.
Soit $(U, R, s, t, c)$ un schéma en groupoïdes sur $S$.
Soit $g : U' \to U$ un morphisme.
Supposons que
\begin{enumerate}
\item le composé
$$
\xymatrix{
U' \times_{g, U, t} R \ar[r]_-{\text{pr}_1} \ar@/^3ex/[rr]^h
& R \ar[r]_s & U
}
$$
a pour image un ouvert $W \subset U$ ;
\item le morphisme induit $h : U' \times_{g, U, t} R \to W$
définit un épimorphisme de faisceaux pour la topologie fppf.
\end{enumerate}
Soit $R' = R|_{U'}$ la restriction de $R$ à $U'$. Alors le morphisme
de faisceaux quotients
$$
U'/R' \to U/R
$$
pour la topologie fppf est représentable et est une immersion ouverte.
\end{lemma}

\begin{proof}
Notons que $W$ est un sous-schéma ouvert $R$-invariant de $U$.
En effet, l'ensemble des points de $W$ est l'ensemble
des points de $U$ équivalents, au sens de
Groupoïdes,
lemme \ref{groupoids-lemma-pre-equivalence-equivalence-relation-points},
à un point de $g(U') \subset U$ (le lemme s'applique puisque $j : R \to U \times_S U$
est une pré-relation d'équivalence d'après
Groupoïdes, lemme \ref{groupoids-lemma-groupoid-pre-equivalence}).
De plus, $g : U' \to U$ se factorise par $W$.
Soit $R|_W$ la restriction de $R$ à $W$.
Il s'ensuit que $R'$ est aussi la restriction de $R|_W$ à $U'$.
On peut donc factoriser le morphisme de faisceaux du lemme sous la forme
$$
U'/R' \longrightarrow W/R|_W \longrightarrow U/R
$$
D'après Groupoïdes, lemme \ref{groupoids-lemma-quotient-groupoid-restrict},
la première flèche est un isomorphisme de faisceaux.
Il suffit donc de démontrer le lemme lorsque $g$ est l'immersion
d'un ouvert $R$-invariant dans $U$.

\medskip\noindent
Supposons que $U' \subset U$ soit un ouvert $R$-invariant et que $g$ soit le
morphisme d'inclusion. Posons $F = U/R$ et $F' = U'/R'$. D'après
Groupoïdes,
lemme \ref{groupoids-lemma-quotient-pre-equivalence-relation-restrict}
ou \ref{groupoids-lemma-quotient-groupoid-restrict},
le morphisme $F' \to F$ est injectif. Soit $\xi \in F(T)$.
Il faut montrer que $T \times_{\xi, F} F'$ est représentable
par un sous-schéma ouvert de $T$.
Il existe un recouvrement fppf $\{f_i : T_i \to T\}$ tel que
$\xi|_{T_i}$ soit l'image par $U \to U/R$ d'un morphisme $a_i : T_i \to U$.
Posons $V_i = a_i^{-1}(U')$.
Affirmons que $V_i \times_T T_j = T_i \times_T V_j$ comme sous-schémas ouverts
de $T_i \times_T T_j$.

\medskip\noindent
Comme $a_i \circ \text{pr}_0$ et $a_j \circ \text{pr}_1$ sont des morphismes
$T_i \times_T T_j \to U$ ayant tous deux pour image la section
$\xi|_{T_i \times_T T_j} \in F(T_i \times_T T_j)$, il existe
un recouvrement fppf $\{f_{ijk} : T_{ijk} \to T_i \times_T T_j\}$ et des morphismes
$r_{ijk} : T_{ijk} \to R$ tels que
$$
a_i \circ \text{pr}_0 \circ f_{ijk} = s \circ r_{ijk},
\quad
a_j \circ \text{pr}_1 \circ f_{ijk} = t \circ r_{ijk},
$$
voir
Groupoïdes, lemme \ref{groupoids-lemma-quotient-pre-equivalence}.
Puisque $U'$ est $R$-invariant, on a $s^{-1}(U') = t^{-1}(U')$, donc
$f_{ijk}^{-1}(V_i \times_T T_j) = f_{ijk}^{-1}(T_i \times_T V_j)$.
Comme $\{f_{ijk}\}$ est surjectif, l'affirmation précédente en résulte.
Ainsi, d'après
Descente, lemme \ref{descent-lemma-open-fpqc-covering},
il existe un sous-schéma ouvert $V \subset T$ tel que
$f_i^{-1}(V) = V_i$. Affirmons que $V$ représente $T \times_{\xi, F} F'$.

\medskip\noindent
Dans un premier temps, montrons que $\xi|_V$ appartient à $F'(V) \subset F(V)$.
En effet, la famille de morphismes $\{V_i \to V\}$ est un recouvrement fppf
et, par construction, $\xi|_{V_i} \in F'(V_i)$.
La propriété de faisceau de $F'$ donne donc $\xi|_V \in F'(V)$.
Enfin, soient $T' \to T$ un morphisme de schémas et
$\xi|_{T'} \in F'(T')$. Pour achever la démonstration, il faut montrer que
$T' \to T$ se factorise par $V$.
Il existe un recouvrement fppf $\{T'_j \to T'\}_{j \in J}$ et des morphismes
$b_j : T'_j \to U'$ tels que $\xi|_{T'_j}$ soit l'image par
$U' \to U/R$ de $b_j$. Il suffit manifestement de montrer que les composés
$T'_j \to T$ se factorisent par $V$. Nous pouvons donc supposer que $\xi|_{T'}$
est l'image d'un morphisme $b : T' \to U'$. Il suffit alors de montrer
que $T'\times_T T_i \to T_i$ se factorise par $V_i$. Sur le schéma
$T' \times_T T_i$, la restriction de $\xi$ est l'image de deux
éléments de $(U/R)(T' \times_T T_i)$, à savoir $a_i \circ \text{pr}_1$ et
$b \circ \text{pr}_0$, dont le second se factorise par l'ouvert $R$-invariant
$U'$. Ainsi, d'après
Groupoïdes, lemme \ref{groupoids-lemma-quotient-pre-equivalence},
il existe un recouvrement $\{h_k : Z_k \to T' \times_T T_i\}$ et des morphismes
$r_k : Z_k \to R$ tels que $a_i \circ \text{pr}_1 \circ h_k = s \circ r_k$
et $b \circ \text{pr}_0 \circ h_k = t \circ r_k$. Comme $U'$ est un ouvert
$R$-invariant et que l'image de $b$ est contenue dans $U'$, l'image de chaque
$a_i \circ \text{pr}_1 \circ h_k$ est contenue dans $U'$. Il s'ensuit que
l'image de $T' \times_T T_i \to T_i$ est contenue dans $V_i$, par définition de $V_i$,
ce qui achève la démonstration.
\end{proof}












\section{Découpage des relations d'équivalence}
\label{section-slicing}

\noindent
Dans cette section, nous expliquons comment « améliorer » une relation
d'équivalence donnée en la découpant. Il ne s'agit pas du type de « découpage
étale » auquel on pourrait être habitué, mais d'un découpage beaucoup plus grossier.


\begin{lemma}
\label{lemma-slice-equivalence-relation}
Soit $S$ un schéma.
Soit $j : R \to U \times_S U$ une relation d'équivalence sur des schémas au-dessus de $S$.
Supposons que $s, t : R \to U$ soient plats et localement de présentation finie.
Alors il existe une relation d'équivalence $j' : R' \to U'\times_S U'$
sur des schémas au-dessus de $S$, et un isomorphisme
$$
U'/R' \longrightarrow U/R
$$
induit par un morphisme $U' \to U$ qui envoie $R'$ dans $R$, tel que
$s', t' : R \to U$ soient plats, localement de présentation finie
et localement quasi-finis.
\end{lemma}

\begin{proof}
Nous démontrerons ce lemme en plusieurs étapes. Nous utiliserons sans plus
le mentionner le fait qu'une relation d'équivalence donne lieu à un schéma en
groupoïdes et que la restriction d'une relation d'équivalence est encore une relation
d'équivalence, voir
Groupoïdes, lemmes
\ref{groupoids-lemma-restrict-relation},
\ref{groupoids-lemma-equivalence-groupoid}, and
\ref{groupoids-lemma-restrict-groupoid-relation}.

\medskip\noindent
Étape 1 : nous pouvons supposer que $s, t : R \to U$ sont des morphismes
localement de présentation finie et de Cohen--Macaulay. En effet, comme dans
Compléments sur les groupoïdes, lemme \ref{more-groupoids-lemma-make-CM},
soit $g : U' \to U$ le sous-schéma ouvert tel que
$t^{-1}(U') \subset R$ soit l'ouvert maximal au-dessus duquel $s : R \to U$ est
de Cohen--Macaulay, et notons $R'$ la restriction de $R$ à $U'$.
D'après le lemme cité, le morphisme
$$
\xymatrix{
t^{-1}(U') \ar@{=}[r] &
U' \times_{g, U, t} R \ar[r]_-{\text{pr}_1} \ar@/^3ex/[rr]^h &
R \ar[r]_s &
U
}
$$
est surjectif. Comme $h$ est plat et localement de présentation finie,
$\{h\}$ est un recouvrement fppf. Ainsi, d'après
Groupoïdes, lemme \ref{groupoids-lemma-quotient-groupoid-restrict},
$U'/R' \to U/R$ est un isomorphisme. Par construction de $U'$,
$s', t'$ sont de Cohen--Macaulay et localement de présentation finie.

\medskip\noindent
Étape 2. Supposons $s, t$ de Cohen--Macaulay et localement de présentation finie.
Soit $u \in U$ un point de type fini. D'après
Compléments sur les groupoïdes, lemme \ref{more-groupoids-lemma-max-slice-quasi-finite},
il existe un schéma affine $U'$ et un morphisme $g : U' \to U$ tels que
\begin{enumerate}
\item $g$ est une immersion ;
\item $u \in U'$ ;
\item $g$ est localement de présentation finie ;
\item $h$ est plat, localement de présentation finie et localement quasi-fini ;
\item les morphismes $s', t' : R' \to U'$ sont plats, localement de
présentation finie et localement quasi-finis.
\end{enumerate}
Nous avons utilisé ici les notations introduites dans
Compléments sur les groupoïdes, situation \ref{more-groupoids-situation-slice}.

\medskip\noindent
Étape 3. Pour chaque point $u \in U$ de type fini,
choisissons $g_u : U'_u \to U$ comme à
l'étape 2 et notons $R'_u$ la restriction de $R$ à $U'_u$.
Notons $h_u = s \circ \text{pr}_1 : U'_u \times_{g_u, U, t} R \to U$. Posons
$U' = \coprod_{u \in U} U'_u$ et $g = \coprod g_u$. Soit $R'$ la
restriction de $R$ à $U'$ comme ci-dessus. Affirmons que
le couple $(U', g)$ convient\footnote{Il faut ici vérifier que $U'$ n'est pas
trop grand, c'est-à-dire qu'il est isomorphe à un objet de la catégorie
$\Sch_{fppf}$, voir la
section \ref{section-conventions}.
Il s'agit d'une question purement ensembliste ; utilisons la notion de taille d'un
schéma introduite dans
Ensembles, section \ref{sets-section-categories-schemes}.
Notons que la taille de chaque $U'_u$ est au plus celle de $U$
et que le cardinal de l'ensemble d'indices est au plus le cardinal de
$|U|$, lui-même borné par la taille de $U$. Ainsi, $U'$ est isomorphe
à un objet de $\Sch_{fppf}$ d'après
Ensembles, lemme \ref{sets-lemma-what-is-in-it}, partie (6).}.
Notons que
\begin{align*}
R' = &
\coprod\nolimits_{u_1, u_2 \in U}
(U'_{u_1} \times_{g_{u_1}, U, t} R)
\times_R
(R \times_{s, U, g_{u_2}} U'_{u_2}) \\
= &
\coprod\nolimits_{u_1, u_2 \in U}
(U'_{u_1} \times_{g_{u_1}, U, t} R) \times_{h_{u_1}, U, g_{u_2}} U'_{u_2}
\end{align*}
Ainsi, la projection $s' : R' \to U' = \coprod U'_{u_2}$
est plate, localement de
présentation finie et localement quasi-finie, car elle est un changement de base de $\coprod h_{u_1}$.
Enfin, par construction, le morphisme
$h : U' \times_{g, U, t} R \to U$ est égal à $\coprod h_u$ ; son
image contient donc tous les points de type fini de $U$.
Comme chaque $h_u$ est plat et localement de présentation finie, on en déduit que
$h$ est plat et localement de présentation finie.
En particulier, l'image de $h$ est ouverte (voir
Morphismes, lemme \ref{morphisms-lemma-fppf-open})
et, puisque l'ensemble des points de type fini est dense (voir
Morphismes, lemme \ref{morphisms-lemma-enough-finite-type-points}),
on en conclut que l'image de $h$ est $U$. Il s'ensuit que
$\{h\}$ est un recouvrement fppf. D'après
Groupoïdes, lemme \ref{groupoids-lemma-quotient-groupoid-restrict},
cela signifie que $U'/R' \to U/R$ est un isomorphisme.
La démonstration du lemme est achevée.
\end{proof}










\section{Quotient par un sous-groupoïde}
\label{section-dividing}

\noindent
Il nous faut encore un lemme avant de pouvoir effectuer l'amorçage final.
Examinons d'abord ce qui se passe pour les groupoïdes « ordinaires », avant
d'aborder la version schématique.

\medskip\noindent
Soit $\mathcal{C}$ un groupoïde, voir
Catégories, définition \ref{categories-definition-groupoid}.
Comme expliqué dans
Groupoïdes, section \ref{groupoids-section-groupoids},
il correspond à un quintuple $(\text{Ob}, \text{Arrows}, s, t, c)$.
Supposons donné un sous-ensemble $P \subset \text{Arrows}$ tel que
$(\text{Ob}, P, s|_P, t|_P, c|_P)$ soit lui aussi un groupoïde et que
$P$ ne contienne aucun automorphisme non trivial. On peut alors construire
le groupoïde quotient
$(\overline{\text{Ob}}, \overline{\text{Arrows}}, \overline{s},
\overline{t}, \overline{c})$
de la manière suivante :
\begin{enumerate}
\item $\overline{\text{Ob}} = \text{Ob}/P$
est l'ensemble des classes de $P$-isomorphie ;
\item $\overline{\text{Arrows}} = P\backslash \text{Arrows}/P$
est l'ensemble des flèches de $\mathcal{C}$ modulo précomposition et
postcomposition par des flèches de $P$ ;
\item les applications source et but
$\overline{s}, \overline{t} : P\backslash \text{Arrows}/P \to \text{Ob}/P$
sont induites par $s, t$ ;
\item la composition est définie par la règle
$\overline{c}(\overline{a}, \overline{b}) = \overline{c(a, b)}$,
qui est bien définie.
\end{enumerate}
En fait, le groupoïde initial
$(\text{Ob}, \text{Arrows}, s, t, c)$ est canoniquement
isomorphe à la restriction (voir la discussion dans
Groupoïdes, section \ref{groupoids-section-restrict-groupoid})
du groupoïde
$(\overline{\text{Ob}}, \overline{\text{Arrows}}, \overline{s},
\overline{t}, \overline{c})$ par l'application quotient
$g : \text{Ob} \to \overline{\text{Ob}}$. Rappelons que cela signifie
que
$$
\text{Arrows} =
\text{Ob}
\times_{g, \overline{\text{Ob}}, \overline{t}}
\overline{\text{Arrows}}
\times_{\overline{s}, \overline{\text{Ob}}, g}
\text{Ob}
$$
ce qui est vrai puisque $P$ ne possède aucun automorphisme non trivial.
Nous omettons les détails.

\medskip\noindent
Le lemme suivant vaut dans une généralité bien plus grande, mais c'est
la version que nous utiliserons dans la démonstration de l'amorçage final (après quoi
nous pourrons démontrer plus facilement les versions plus générales de ce lemme).

\begin{lemma}
\label{lemma-divide-subgroupoid}
Soit $S$ un schéma.
Soit $(U, R, s, t, c)$ un schéma en groupoïdes sur $S$.
Soit $P \to R$ un monomorphisme de schémas. Supposons que
\begin{enumerate}
\item $(U, P, s|_P, t|_P, c|_{P \times_{s, U, t}P})$ soit un schéma en groupoïdes ;
\item $s|_P, t|_P : P \to U$ soient finis localement libres ;
\item $j|_P : P \to U \times_S U$ soit un monomorphisme ;
\item $U$ soit affine ;
\item $j : R \to U \times_S U$ soit séparé et localement quasi-fini.
\end{enumerate}
Alors $U/P$ est représentable par un schéma affine $\overline{U}$, le
morphisme quotient $U \to \overline{U}$ est fini localement libre et
$P = U \times_{\overline{U}} U$. De plus, $R$ est la restriction d'un
schéma en groupoïdes
$(\overline{U}, \overline{R}, \overline{s}, \overline{t}, \overline{c})$
sur $\overline{U}$ par le morphisme quotient $U \to \overline{U}$.
\end{lemma}

\begin{proof}
Les conditions (1), (2), (3) et (4), ainsi que
Groupoïdes, proposition \ref{groupoids-proposition-finite-flat-equivalence},
impliquent l'existence du schéma affine $\overline{U}$ représentant $U/P$ ;
le morphisme $U \to \overline{U}$ est fini localement libre et
$P = U \times_{\overline{U}} U$. L'identification
$P = U \times_{\overline{U}} U$ est telle que $t|_P = \text{pr}_0$ et
$s|_P = \text{pr}_1$, et que la composition est égale à
$\text{pr}_{02} : U \times_{\overline{U}} U \times_{\overline{U}} U
\to U \times_{\overline{U}} U$.
Un produit de morphismes finis localement libres est fini localement libre (voir
Espaces, lemme \ref{spaces-lemma-product-representable-transformations-property}
et
Morphismes, lemmes \ref{morphisms-lemma-base-change-finite-locally-free} et
\ref{morphisms-lemma-composition-finite-locally-free}).
Pour obtenir $\overline{R}$, nous allons descendre
le schéma $R$ par le morphisme fini localement libre
$U \times_S U \to \overline{U} \times_S \overline{U}$.
Notons en effet que
$$
(U \times_S U)
\times_{(\overline{U} \times_S \overline{U})}
(U \times_S U)
=
P \times_S P
$$
d'après ce qui précède. Ainsi, se donner une donnée de descente (voir
Descente, définition \ref{descent-definition-descent-datum})
pour $R / U \times_S U / \overline{U} \times_S \overline{U}$
revient à se donner un isomorphisme
$$
\varphi :
R \times_{(U \times_S U), t \times t} (P \times_S P)
\longrightarrow
(P \times_S P) \times_{s \times s, (U \times_S U)} R
$$
au-dessus de $P \times_S P$ satisfaisant à une condition de cocycle. Nous définissons $\varphi$
sur les points à valeurs dans $T$ par la règle
$$
\varphi : (r, (p, p')) \longmapsto ((p, p'), p^{-1} \circ r \circ p')
$$
où la composition est prise dans la catégorie groupoïde
$(U(T), R(T), s, t, c)$.
Cela a un sens, car pour que $(r, (p, p'))$ soit un point à valeurs dans $T$
de la source de $\varphi$, il faut que $t(r) = t(p)$
et $s(r) = t(p')$. Notons que ce morphisme est un isomorphisme
dont l'inverse est donné par
$((p, p'), r') \mapsto (p \circ r' \circ (p')^{-1}, (p, p'))$.
Pour vérifier la condition de cocycle, il faut établir que
$\varphi_{02} = \varphi_{12} \circ \varphi_{01}$
comme morphismes au-dessus de
$$
(U \times_S U)
\times_{(\overline{U} \times_S \overline{U})} (U \times_S U)
\times_{(\overline{U} \times_S \overline{U})} (U \times_S U) =
(P \times_S P) \times_{s \times s, (U \times_S U), t \times t} (P \times_S P)
$$
Un calcul explicite donne
$$
\begin{matrix}
\varphi_{02} & (r, (p_1, p_1'), (p_2, p_2')) & \mapsto &
((p_1, p_1'), (p_2, p_2'),
(p_1 \circ p_2)^{-1} \circ r \circ (p_1' \circ p_2')) \\
\varphi_{01} & (r, (p_1, p_1'), (p_2, p_2')) & \mapsto &
((p_1, p_1'), p_1^{-1} \circ r \circ p_1', (p_2, p_2')) \\
\varphi_{12} & ((p_1, p_1'), r, (p_2, p_2')) & \mapsto &
((p_1, p_1'), (p_2, p_2'), p_2^{-1} \circ r \circ p_2')
\end{matrix}
$$
(avec des notations évidentes), ce qui donne le résultat voulu.
Comme $j$ est séparé et localement quasi-fini d'après (5), nous pouvons appliquer
Compléments sur les morphismes, lemme
\ref{more-morphisms-lemma-separated-locally-quasi-finite-morphisms-fppf-descend}
pour obtenir un schéma $\overline{R} \to \overline{U} \times_S \overline{U}$
et un isomorphisme
$$
R \to \overline{R} \times_{(\overline{U} \times_S \overline{U})} (U \times_S U)
$$
qui identifie la donnée de descente $\varphi$ à la donnée de descente
canonique sur
$\overline{R} \times_{(\overline{U} \times_S \overline{U})} (U \times_S U)$,
voir
Descente, définition \ref{descent-definition-effective}.

\medskip\noindent
Comme $U \times_S U \to \overline{U} \times_S \overline{U}$ est fini
localement libre, on en déduit que $R \to \overline{R}$ est fini localement libre
comme changement de base. Ainsi, $R \to \overline{R}$ est surjectif comme morphisme de
faisceaux sur $(\Sch/S)_{fppf}$.
Notre choix de $\varphi$ implique que, pour des points à valeurs dans $T$, $r, r' \in R(T)$,
ceux-ci ont la même image dans $\overline{R}$ si et seulement si
$p^{-1} \circ r \circ p'$ pour certains $p, p' \in P(T)$. Ainsi,
$\overline{R}$ représente le faisceau
$$
T \longmapsto  \overline{R(T)} = P(T)\backslash R(T)/P(T)
$$
avec les notations de la discussion précédant le lemme.
Nous pouvons donc définir la structure de groupoïde sur
$(\overline{U} = U/P, \overline{R} = P\backslash R/P)$ exactement comme dans
la discussion du cas d'un groupoïde « ordinaire ».
Il s'ensuit que $(U, R, s, t, c)$ est l'image inverse de
cette structure de groupoïde par le morphisme $U \to \overline{U}$.
Cela achève la démonstration.
\end{proof}















\section{Amorçage final}
\label{section-final-bootstrap}

\noindent
Le résultat suivant va nettement plus loin que les précédents.

\begin{theorem}
\label{theorem-final-bootstrap}
Soit $S$ un schéma.
Soit $F : (\Sch/S)_{fppf}^{opp} \to \textit{Sets}$ un foncteur.
Chacune des conditions suivantes implique que $F$ est un espace algébrique :
\begin{enumerate}
\item $F = U/R$, où $(U, R, s, t, c)$ est un groupoïde en espaces algébriques
sur $S$ tel que $s, t$ soient plats et localement de présentation finie, et
$j = (t, s) : R \to U \times_S U$ soit une relation d'équivalence ;
\item $F = U/R$, où $(U, R, s, t, c)$ est un schéma en groupoïdes
sur $S$ tel que $s, t$ soient plats et localement de présentation finie, et
$j = (t, s) : R \to U \times_S U$ soit une relation d'équivalence ;
\item $F$ est un faisceau et il existe un espace algébrique $U$ et un morphisme
$U \to F$ représentable par des espaces algébriques,
surjectif, plat et localement de présentation finie ;
\item $F$ est un faisceau et il existe un schéma $U$ et un morphisme
$U \to F$ représentable par des espaces algébriques ou par des schémas,
surjectif, plat et localement de présentation finie ;
\item $F$ est un faisceau, $\Delta_F$ est représentable par des espaces algébriques,
et il existe un espace algébrique $U$ et un morphisme $U \to F$
surjectif, plat et localement de présentation finie ;
\item $F$ est un faisceau, $\Delta_F$ est représentable,
et il existe un schéma $U$ et un morphisme $U \to F$
surjectif, plat et localement de présentation finie.
\end{enumerate}
\end{theorem}

\begin{proof}
Observations immédiates : (6) est un cas particulier de (5) et
(4) est un cas particulier de (3).
Montrons d'abord que les cas (5) et (3) se ramènent au cas (1).
En effet, le lemme d'amorçage de la diagonale
\ref{lemma-bootstrap-diagonal}
montre que (3) implique (5). Dans le cas (5), posons $R = U \times_F U$ ; c'est
un espace algébrique par hypothèse. De plus, toujours par hypothèse, les deux
projections $s, t : R \to U$ sont surjectives, plates et localement de
présentation finie. Le morphisme $j : R \to U \times_S U$ est manifestement une
relation d'équivalence. D'après le
lemme \ref{lemma-surjective-flat-locally-finite-presentation},
le morphisme $U \to F$ est un épimorphisme de faisceaux. Ainsi $F = U/R$,
ce qui nous ramène au cas (1).

\medskip\noindent
Montrons ensuite que (1) se ramène à (2).
Soit $(U, R, s, t, c)$ un groupoïde en espaces algébriques
sur $S$ tel que $s, t$ soient plats et localement de présentation finie, et que
$j = (t, s) : R \to U \times_S U$ soit une relation d'équivalence.
Choisissons un schéma $U'$ et un morphisme étale surjectif $U' \to U$.
Soit $R' = R|_{U'}$ la restriction de $R$ à $U'$. D'après
Groupoïdes dans les espaces,
lemme \ref{spaces-groupoids-lemma-quotient-pre-equivalence-relation-restrict},
on a $U/R = U'/R'$. Comme $s', t' : R' \to U'$ sont eux aussi
plats et localement de présentation finie (voir
Compléments sur les groupoïdes dans les espaces,
lemme \ref{spaces-more-groupoids-lemma-restrict-preserves-type}),
nous sommes ramenés au cas où $U$ est un schéma.
Puisque $j$ est une relation d'équivalence, $j$ est un monomorphisme.
Comme $s : R \to U$ est localement de présentation finie,
$j : R \to U \times_S U$ est localement de type fini, voir
Morphismes d'espaces, lemme \ref{spaces-morphisms-lemma-permanence-finite-type}.
D'après
Morphismes d'espaces, lemme
\ref{spaces-morphisms-lemma-monomorphism-loc-finite-type-loc-quasi-finite}
$j$ est localement quasi-fini et séparé.
Ainsi, si $U$ est un schéma, alors $R$ est un schéma d'après
Morphismes d'espaces, proposition
\ref{spaces-morphisms-proposition-locally-quasi-finite-separated-over-scheme}.
Nous sommes donc ramenés à démontrer le théorème dans le cas (2).

\medskip\noindent
Supposons $F = U/R$, où $(U, R, s, t, c)$ est un schéma en groupoïdes
sur $S$ tel que $s, t$ soient plats et localement de présentation finie, et que
$j = (t, s) : R \to U \times_S U$ soit une relation d'équivalence. D'après le
lemme \ref{lemma-slice-equivalence-relation},
nous sommes ramenés au cas où $s, t$ sont plats,
localement de présentation finie et localement quasi-finis.
Soit $U = \bigcup_{i \in I} U_i$ un recouvrement ouvert affine
(où l'ensemble d'indices $I$ est de cardinal $\leq$ à la taille de $U$, afin d'éviter
plus loin des problèmes ensemblistes --- la plupart des lecteurs peuvent ignorer cette remarque).
Soit $(U_i, R_i, s_i, t_i, c_i)$ la restriction de $R$
à $U_i$. Il est clair que $s_i, t_i$ restent plats, localement de présentation
finie et localement quasi-finis, puisque $R_i$ est le sous-schéma ouvert
$s^{-1}(U_i) \cap t^{-1}(U_i)$ de $R$
et que $s_i, t_i$ sont les restrictions de $s, t$ à cet ouvert. D'après le
lemme \ref{lemma-better-finding-opens}
(ou, plus simplement,
Espaces, lemme \ref{spaces-lemma-finding-opens}),
le morphisme $U_i/R_i \to U/R$ est représentable par des immersions ouvertes.
Ainsi, si l'on peut montrer que $F_i = U_i/R_i$ est un espace algébrique, alors
$\coprod_{i \in I} F_i$ est un espace algébrique d'après
Espaces, lemme \ref{spaces-lemma-coproduct-algebraic-spaces}.
Comme $U = \bigcup U_i$ est un recouvrement ouvert, il est clair que
$\coprod F_i \to F$ est surjectif. Il s'ensuit que
$U/R$ est un espace algébrique, d'après
Espaces, lemme \ref{spaces-lemma-glueing-algebraic-spaces}.
Nous sommes ainsi ramenés au cas où $U$ est affine, $s, t$ sont plats,
localement de présentation finie et localement quasi-finis, et où
$j$ est une relation d'équivalence.

\medskip\noindent
Supposons que $(U, R, s, t, c)$ soit un schéma en groupoïdes sur $S$,
avec $U$ affine, tel que $s, t$ soient plats, localement de présentation finie
et localement quasi-finis, et que $j$ soit une relation d'équivalence.
Choisissons $u \in U$. Appliquons
Compléments sur les groupoïdes dans les espaces,
lemme \ref{spaces-more-groupoids-lemma-quasi-splitting-affine-scheme},
à $u \in U, R, s, t, c$. On obtient un schéma affine $U'$, un morphisme étale
$g : U' \to U$ et un point $u' \in U'$ tel que $\kappa(u) = \kappa(u')$
et que la restriction $R' = R|_{U'}$ soit quasi-scindée au-dessus de $u'$.
Notons que l'image $g(U')$ est ouverte, puisque $g$ est étale, et qu'elle contient $u$.
En appliquant le lemme à plusieurs reprises, on peut donc trouver un nombre fini de
points $u_i \in U$, $i = 1, \ldots, n$,
de schémas affines $U'_i$, de morphismes étales $g_i : U_i' \to U$ et de points
$u'_i \in U'_i$ tels que $g(u'_i) = u_i$, avec (a) chaque
restriction $R'_i$ quasi-scindée au-dessus d'un point de $U'_i$, et
(b) $U = \bigcup_{i = 1, \ldots, n} g_i(U'_i)$.
Reprenons maintenant la dernière partie de l'argument de l'alinéa précédent :
d'après le
lemme \ref{lemma-better-finding-opens}
(ou, plus simplement,
Espaces, lemme \ref{spaces-lemma-finding-opens}),
le morphisme $U'_i/R'_i \to U/R$ est représentable par des immersions ouvertes.
Si l'on peut montrer que $F_i = U'_i/R'_i$ est un espace algébrique, alors
$\coprod_{i \in I} F_i$ est un espace algébrique d'après
Espaces, lemme \ref{spaces-lemma-coproduct-algebraic-spaces}.
Comme $\{g_i : U'_i \to U\}$ est un recouvrement étale,
il est clair que $\coprod F_i \to F$ est surjectif. Il s'ensuit que
$U/R$ est un espace algébrique, d'après
Espaces, lemme \ref{spaces-lemma-glueing-algebraic-spaces}.
Nous sommes ainsi ramenés au cas où $U$ est affine, $s, t$ sont plats,
localement de présentation finie et localement quasi-finis,
$j$ est une relation d'équivalence et $R$ est quasi-scindé au-dessus de $u$ pour un certain
$u \in U$.

\medskip\noindent
Supposons que $(U, R, s, t, c)$ soit un schéma en groupoïdes sur $S$,
avec $U$ affine et $u \in U$, tel que $s, t$ soient plats, localement
de présentation finie et localement quasi-finis, que
$j = (t, s) : R \to U \times_S U$ soit une relation d'équivalence
et que $R$ soit quasi-scindé au-dessus de $u$. Soit $P \subset R$ un quasi-scindage
de $R$ au-dessus de $u$. D'après le
lemme \ref{lemma-divide-subgroupoid},
$(U, R, s, t, c)$ est la restriction d'un groupoïde
$(\overline{U}, \overline{R}, \overline{s}, \overline{t}, \overline{c})$
par un morphisme fini localement libre surjectif $U \to \overline{U}$ tel que
$P = U \times_{\overline{U}} U$. Notons que $s$ admet la factorisation
$$
R = U \times_{\overline{U}, \overline{t}} \overline{R}
\times_{\overline{s}, \overline{U}} U
\xrightarrow{\text{pr}_{23}}
\overline{R} \times_{\overline{s}, \overline{U}} U
\xrightarrow{\text{pr}_2} U
$$
Le morphisme $\text{pr}_2$ est le changement de base de $\overline{s}$, et
$\text{pr}_{23}$ est un changement de base du morphisme fini localement libre
surjectif $U \to \overline{U}$. Comme $s$ est plat, localement
de présentation finie et localement quasi-fini, et comme $\text{pr}_{23}$
est fini localement libre surjectif (comme changement de base d'un tel morphisme),
on en déduit que $\text{pr}_2$ est plat, localement
de présentation finie et localement quasi-fini d'après
Descente, lemmes
\ref{descent-lemma-flat-fpqc-local-source} et
\ref{descent-lemma-locally-finite-presentation-fppf-local-source}, ainsi que
Morphismes, lemme \ref{morphisms-lemma-quasi-finite-local-source}.
Comme $\text{pr}_2$ est le changement de base du morphisme
$\overline{s}$ par $U \to \overline{U}$ et que $\{U \to \overline{U}\}$
est un recouvrement fppf, on en conclut que $\overline{s}$ est
plat, localement de présentation finie et localement quasi-fini, voir
Descente, lemmes \ref{descent-lemma-descending-property-flat},
\ref{descent-lemma-descending-property-locally-finite-presentation} et
\ref{descent-lemma-descending-property-quasi-finite}. Il en va de même
pour $\overline{t}$. Considérons le diagramme commutatif
$$
\xymatrix{
U \times_{\overline{U}} U \ar@{=}[r] \ar[rd] & P \ar[r] \ar[d] & R \ar[d] \\
& \overline{U} \ar[r]^{\overline{e}} & \overline{R}
}
$$
Un fait général sur les restrictions affirme que les quatre sommets extérieurs
forment un diagramme cartésien. L'égalité montre que le carré intérieur est
cartésien. Comme $P$ est ouvert dans $R$ (par définition d'un quasi-scindage),
on en conclut que $\overline{e}$ est une immersion ouverte d'après
Descente, lemme \ref{descent-lemma-descending-property-open-immersion}.
En appliquant
Groupoïdes,
lemme \ref{groupoids-lemma-quotient-pre-equivalence-relation-restrict},
on obtient $U/R = \overline{U}/\overline{R}$. Nous sommes donc ramenés
au cas où $(U, R, s, t, c)$ est un schéma en groupoïdes sur $S$,
avec $U$ affine et $u \in U$, tel que $s, t$ soient plats, localement
de présentation finie et localement quasi-finis, que
$j = (t, s) : R \to U \times_S U$ soit une relation d'équivalence
et que $e : U \to R$ soit une immersion ouverte !

\medskip\noindent
Mais bien entendu, si $e$ est une immersion ouverte et si
$s, t$ sont plats et localement de présentation finie,
alors les morphismes $t, s$ sont étales.
On peut par exemple le voir en appliquant
Compléments sur les groupoïdes, lemme \ref{more-groupoids-lemma-sheaf-differentials},
qui montre que $\Omega_{R/U} = 0$, ce qui implique à son tour
que $s, t : R \to U$ sont G-non ramifiés (voir
Morphismes, lemme \ref{morphisms-lemma-unramified-omega-zero}),
et implique donc que $s, t$ sont étales (voir
Morphismes, lemme \ref{morphisms-lemma-flat-unramified-etale}).
Enfin, si $s, t$ sont étales, alors $U/R$ est un espace
algébrique d'après
Espaces, théorème \ref{spaces-theorem-presentation}.
\end{proof}





\section{Applications}
\label{section-applications}

\noindent
Comme première application, nous obtenons le fait fondamental suivant :
$$
\fbox{Un faisceau qui est localement un espace algébrique pour la topologie fppf est un espace algébrique.}
$$
C'est le contenu du lemme suivant.
Notons que l'hypothèse (2) équivaut à la condition selon laquelle
$F|_{(\Sch/S_i)_{fppf}}$ est un espace algébrique, voir
Espaces, lemme \ref{spaces-lemma-rephrase}.
L'hypothèse (3) est une condition ensembliste que peuvent ignorer
ceux qui ne se soucient pas des questions ensemblistes.

\begin{lemma}
\label{lemma-locally-algebraic-space}
\begin{slogan}
La définition d'un espace algébrique est locale pour la topologie fppf.
\end{slogan}
Soit $S$ un schéma.
Soit $F : (\Sch/S)_{fppf}^{opp} \to \textit{Sets}$ un foncteur.
Soit $\{S_i \to S\}_{i \in I}$ un recouvrement de $(\Sch/S)_{fppf}$.
Supposons que
\begin{enumerate}
\item $F$ soit un faisceau ;
\item chaque $F_i = h_{S_i} \times F$ soit un espace algébrique ;
\item $\coprod_{i \in I} F_i$ soit un espace algébrique (voir
Espaces, lemme \ref{spaces-lemma-coproduct-algebraic-spaces}).
\end{enumerate}
Alors $F$ est un espace algébrique.
\end{lemma}

\begin{proof}
Considérons le morphisme $\coprod F_i \to F$. C'est le changement de base
de $\coprod S_i \to S$ par $F \to S$. Il est donc représentable,
localement de présentation finie, plat et surjectif, d'après notre définition
d'un recouvrement fppf et le
lemme \ref{lemma-base-change-transformation-property}.
Ainsi, le
théorème \ref{theorem-final-bootstrap}
s'applique et montre que $F$ est un espace algébrique.
\end{proof}

\noindent
Voici un cas particulier du lemme \ref{lemma-locally-algebraic-space}
dans lequel les questions ensemblistes ne se posent pas.

\begin{lemma}
\label{lemma-locally-algebraic-space-finite-type}
Soit $S$ un schéma.
Soit $F : (\Sch/S)_{fppf}^{opp} \to \textit{Sets}$ un foncteur.
Soit $\{S_i \to S\}_{i \in I}$ un recouvrement de $(\Sch/S)_{fppf}$.
Supposons que
\begin{enumerate}
\item $F$ soit un faisceau ;
\item chaque $F_i = h_{S_i} \times F$ soit un espace algébrique ;
\item les morphismes $F_i \to S_i$ soient de type fini.
\end{enumerate}
Alors $F$ est un espace algébrique.
\end{lemma}

\begin{proof}
Nous allons utiliser le
lemme \ref{lemma-locally-algebraic-space}
ci-dessus. Pour cela, nous montrerons que l'hypothèse selon laquelle
$F_i$ est de type fini sur $S_i$ entraîne que la condition ensembliste
du lemme est satisfaite (quitte à raffiner légèrement le
recouvrement donné de $S$).
Nous suggérons au lecteur de passer la suite de la démonstration.

\medskip\noindent
Si $S'_i \to S_i$ est un morphisme de schémas, alors
$$
h_{S'_i} \times F =
h_{S'_i} \times_{h_{S_i}} h_{S_i} \times F =
h_{S'_i} \times_{h_{S_i}} F_i
$$
est un espace algébrique de type fini sur $S'_i$, voir
Espaces, lemme \ref{spaces-lemma-fibre-product-spaces},
et
Morphismes d'espaces,
lemme \ref{spaces-morphisms-lemma-base-change-finite-type}.
Nous pouvons donc raffiner le recouvrement donné. Après ce raffinement, nous pouvons supposer :
(a) chaque $S_i$ est affine ; (b) le cardinal de $I$ est au plus
le cardinal de l'ensemble des points de $S$. (En effet, pour recouvrir
tout $S$, il suffit que chaque point appartienne à l'image de $S_i \to S$
pour un certain $i$.)

\medskip\noindent
Comme chaque $S_i$ est affine et chaque $F_i$ de type fini sur $S_i$,
on en déduit que $F_i$ est quasi-compact. Ainsi, d'après
Propriétés des espaces,
lemme \ref{spaces-properties-lemma-quasi-compact-affine-cover},
il existe un $U_i \in \Ob((\Sch/S)_{fppf})$ affine
et un morphisme étale surjectif $U_i \to F_i$. Le fait que
$F_i \to S_i$ soit localement de type fini implique alors que
$U_i \to S_i$ est localement de type fini, et en particulier que
$U_i \to S$ est localement de type fini. D'après
Ensembles, lemme \ref{sets-lemma-bound-finite-type},
on a $\text{size}(U_i) \leq \text{size}(S)$.
Comme on a aussi $|I| \leq \text{size}(S)$,
$\coprod_{i \in I} U_i$ est isomorphe à un objet de
$(\Sch/S)_{fppf}$ d'après
Ensembles, lemme \ref{sets-lemma-bound-size},
et la construction de $\Sch$. Il s'ensuit que
$\coprod F_i$ est un espace algébrique d'après
Espaces, lemme \ref{spaces-lemma-coproduct-algebraic-spaces},
ce qui conclut.
\end{proof}

\noindent
Comme deuxième application, nous obtenons
$$
\fbox{Toute donnée de descente fppf pour des espaces algébriques est effective.}
$$
Cet énoncé vaut sous réserve de difficultés ensemblistes ; nous donnons
le lemme suivant comme résultat type.

\begin{lemma}
\label{lemma-descend-algebraic-space}
\begin{slogan}
Les données de descente fppf pour les espaces algébriques sont effectives.
\end{slogan}
Soit $S$ un schéma. Soit $\{X_i \to X\}_{i \in I}$ un recouvrement
fppf d'espaces algébriques sur $S$.
\begin{enumerate}
\item Si $I$ est dénombrable\footnote{La restriction de dénombrabilité peut être
ignorée par ceux que les questions ensemblistes ne préoccupent pas. On peut admettre
ici des ensembles d'indices plus grands si l'on peut borner la taille des espaces
algébriques que l'on descend. Voir par exemple le
lemme \ref{lemma-locally-algebraic-space-finite-type}.}, alors toute
donnée de descente d'espaces algébriques relativement à $\{X_i \to X\}$ est effective.
\item Toute donnée de descente $(Y_i, \varphi_{ij})$ relativement à
$\{X_i \to X\}_{i \in I}$ (Descente sur les espaces, définition
\ref{spaces-descent-definition-descent-datum-for-family-of-morphisms})
pour laquelle $Y_i \to X_i$ est de type fini
est effective.
\end{enumerate}
\end{lemma}

\begin{proof}
Démonstration de (1). D'après
Descente sur les espaces, lemme \ref{spaces-descent-lemma-descent-data-sheaves},
cela revient à affirmer qu'un faisceau fppf $F$
muni d'un morphisme $F \to X$ est un espace algébrique dès que
chaque $F \times_X X_i$ est un espace algébrique.
La restriction sur le cardinal de $I$ implique que
les coproduits d'espaces algébriques indexés par $I$ sont des espaces algébriques, voir
Espaces, lemme \ref{spaces-lemma-coproduct-algebraic-spaces},
et
Ensembles, lemme \ref{sets-lemma-what-is-in-it}.
Le morphisme
$$
\coprod F \times_X X_i \longrightarrow F
$$
est représentable par des espaces algébriques (comme changement de base de
$\coprod X_i \to X$, voir le lemme \ref{lemma-base-change-transformation}),
et surjectif, plat et localement de présentation finie
(comme changement de base de $\coprod X_i \to X$, voir le
lemme \ref{lemma-base-change-transformation-property}).
La partie (1) résulte donc du théorème \ref{theorem-final-bootstrap}.

\medskip\noindent
Démonstration de (2). Appliquons d'abord
Descente sur les espaces, lemme \ref{spaces-descent-lemma-descent-data-sheaves},
pour obtenir un faisceau fppf $F$ muni d'un morphisme $F \to X$
tel que $F \times_X X_i = Y_i$ pour tout $i \in I$.
Il s'agit de montrer que $F$ est un espace algébrique.
Choisissons un schéma $U$ et un morphisme étale surjectif $U \to X$.
Alors $F' = U \times_X F \to F$ est représentable, surjectif et étale
comme changement de base de $U \to X$.
D'après le théorème \ref{theorem-final-bootstrap}, il suffit de montrer
que $F' = U \times_X F$ est un espace algébrique.
On peut choisir un recouvrement fppf $\{U_j \to U\}_{j \in J}$,
où $U_j$ est un schéma raffinant le recouvrement fppf
$\{X_i \times_X U \to U\}_{i \in I}$, voir
Topologies sur les espaces, lemme
\ref{spaces-topologies-lemma-refine-fppf-schemes}.
On obtient ainsi une application $a : J \to I$ et, pour chaque $j$,
un morphisme $U_j \to X_{a(j)}$ au-dessus de $X$.
Alors $U_j \times_U F' = U_j \times_{X_{a(j)}} Y_{a(j)}$
est de type fini sur $U_j$. Par conséquent, $F'$ est un espace
algébrique d'après le lemme \ref{lemma-locally-algebraic-space-finite-type}.
\end{proof}

\noindent
Voici une application d'un autre type.

\begin{lemma}
\label{lemma-representable-by-spaces-cover}
Soit $S$ un schéma. Soient $a : F \to G$ et $b : G \to H$ des
transformations de foncteurs $(\Sch/S)_{fppf}^{opp} \to \textit{Sets}$.
Supposons que
\begin{enumerate}
\item $F, G, H$ soient des faisceaux ;
\item $a : F \to G$ soit représentable par des espaces algébriques, plat,
localement de présentation finie et surjectif ;
\item $b \circ a : F \to H$ soit représentable par des espaces algébriques.
\end{enumerate}
Alors $b$ est représentable par des espaces algébriques.
\end{lemma}

\begin{proof}
Soit $U$ un schéma sur $S$ et soit $\xi \in H(U)$. Il faut montrer que
$U \times_{\xi, H} G$ est un espace algébrique. Or, nous savons
que $U \times_{\xi, H} F$ est un espace algébrique et que
$U \times_{\xi, H} F \to U \times_{\xi, H} G$ est représentable par
des espaces algébriques, plat, localement de présentation finie et surjectif,
comme changement de base du morphisme $a$ (voir le
lemme \ref{lemma-base-change-transformation-property}).
Le résultat découle donc du théorème \ref{theorem-final-bootstrap}.
\end{proof}

\begin{lemma}
\label{lemma-quotient-stack-isom}
Supposons que $B \to S$ et $(U, R, s, t, c)$ soient comme dans
Groupoïdes dans les espaces,
définition \ref{spaces-groupoids-definition-quotient-stack} (1).
Pour tout schéma $T$ sur $S$ et tous objets $x, y$ de $[U/R]$ au-dessus de $T$,
le faisceau $\mathit{Isom}(x, y)$ sur $(\Sch/T)_{fppf}$
est un espace algébrique.
\end{lemma}

\begin{proof}
D'après
Groupoïdes dans les espaces,
lemme \ref{spaces-groupoids-lemma-quotient-stack-isom},
il existe un recouvrement fppf $\{T_i \to T\}_{i \in I}$
tel que $\mathit{Isom}(x, y)|_{(\Sch/T_i)_{fppf}}$
soit un espace algébrique pour chaque $i$. D'après
Espaces, lemme \ref{spaces-lemma-rephrase},
cela signifie que chaque $F_i = h_{S_i} \times \mathit{Isom}(x, y)$
est un espace algébrique.
Pour démontrer le lemme, il ne reste donc qu'à vérifier la condition ensembliste
du lemme \ref{lemma-locally-algebraic-space} ci-dessus, à savoir que
$\coprod F_i$ est un espace algébrique. Pour cela, nous utilisons
Espaces, lemme \ref{spaces-lemma-coproduct-algebraic-spaces},
qui demande de montrer que $I$ et les $F_i$ ne sont pas « trop grands ».
Nous suggérons au lecteur de passer la suite de la démonstration.

\medskip\noindent
Choisissons $U' \in \Ob(\Sch/S)_{fppf}$ et un morphisme
étale surjectif $U' \to U$. Soit $R'$ la restriction de $R$ à $U'$.
Comme $[U/R] = [U'/R']$, nous pouvons, après avoir remplacé $U$ par $U'$,
supposer que $U$ est un schéma. (Cette étape sert à faire en sorte que les
produits fibrés ci-dessous soient pris au-dessus d'un schéma.)

\medskip\noindent
Notons que, si l'on raffine le recouvrement $\{T_i \to T\}$, chaque
$F_i$ reste un espace algébrique.
Nous pouvons donc supposer que chaque $T_i$ est affine. Comme
$T_i \to T$ est localement de présentation finie, il s'ensuit que
$\text{size}(T_i) \leq \text{size}(T)$, voir
Ensembles, lemme \ref{sets-lemma-bound-finite-type}.
Nous pouvons aussi supposer que le cardinal de l'ensemble d'indices $I$ est au plus le
cardinal de l'ensemble des points de $T$, car, pour obtenir un
recouvrement, il suffit que chaque point de $T$ appartienne à l'image.
Ainsi $|I| \leq \text{size}(T)$.
Choisissons $W \in \Ob((\Sch/S)_{fppf})$
et un morphisme étale surjectif $W \to R$. Notons que, dans la démonstration de
Groupoïdes dans les espaces,
lemme \ref{spaces-groupoids-lemma-quotient-stack-isom},
nous avons montré que $F_i$ est représenté par
$T_i \times_{(y_i, x_i), U \times_B U} R$ pour certains
$x_i, y_i : T_i \to U$. On voit alors que
$V_i = T_i \times_{(y_i, x_i), U \times_B U} W$ est un
schéma muni d'un morphisme étale surjectif $V_i \to F_i$.
D'après
Ensembles, lemme \ref{sets-lemma-bound-size-fibre-product},
on a
$$
\text{size}(V_i) \leq \max\{\text{size}(T_i), \text{size}(W)\}
\leq \max\{\text{size}(T), \text{size}(W)\}
$$
Ainsi, d'après
Ensembles, lemme \ref{sets-lemma-bound-size},
on en déduit que
$$
\text{size}(\coprod\nolimits_{i \in I} V_i)
\leq \max\{|I|, \text{size}(T), \text{size}(W)\}.
$$
La construction de $\Sch$ montre donc
que $\coprod_{i \in I} V_i$ est isomorphe à un objet
$V$ de $(\Sch/S)_{fppf}$. Cela vérifie
l'hypothèse de
Espaces, lemme \ref{spaces-lemma-coproduct-algebraic-spaces},
ce qui conclut.
\end{proof}

\begin{lemma}
\label{lemma-covering-quotient}
Soit $S$ un schéma. Considérons un espace algébrique $F$ de la forme $F = U/R$,
où $(U, R, s, t, c)$ est un groupoïde en espaces algébriques
sur $S$ tel que $s, t$ soient plats et localement de présentation finie, et que
$j = (t, s) : R \to U \times_S U$ soit une relation d'équivalence.
Alors $U \to F$ est surjectif, plat et localement de présentation finie.
\end{lemma}

\begin{proof}
Ce résultat est presque, mais pas tout à fait, immédiat. En effet, d'après
Groupoïdes dans les espaces, lemme
\ref{spaces-groupoids-lemma-quotient-pre-equivalence}
et puisque $j$ est un monomorphisme, on a $R = U \times_F U$.
Choisissons un schéma $W$ et un morphisme étale surjectif $W \to F$.
Comme $U \to F$ est un épimorphisme de faisceaux, il existe un recouvrement fppf
$\{W_i \to W\}$ et des morphismes $W_i \to U$ relevant les morphismes $W_i \to F$.
Alors
$$
W_i \times_F U = W_i \times_U U \times_F U = W_i \times_{U, t} R
$$
et la projection $W_i \times_F U \to W_i$ est le changement de base de
$t : R \to U$, donc elle est plate et localement de présentation finie, voir
Morphismes d'espaces, lemmes
\ref{spaces-morphisms-lemma-base-change-flat} and
\ref{spaces-morphisms-lemma-base-change-finite-presentation}.
Ainsi, d'après
Descente sur les espaces, lemmes
\ref{spaces-descent-lemma-descending-property-flat} and
\ref{spaces-descent-lemma-descending-property-locally-finite-presentation}
on voit que $U \to F$ est plat et localement de présentation finie.
Il est surjectif d'après
Espaces, remarque \ref{spaces-remark-warning}.
\end{proof}

\begin{lemma}
\label{lemma-quotient-free-action}
Soit $S$ un schéma. Soit $X \to B$ un morphisme d'espaces algébriques sur
$S$. Soit $G$ un espace algébrique en groupes sur $B$ et soit
$a : G \times_B X \to X$ une action de $G$ sur $X$ au-dessus de $B$.
Si
\begin{enumerate}
\item $a$ est une action libre ;
\item $G \to B$ est plat et localement de présentation finie,
\end{enumerate}
alors $X/G$ (voir
Groupoïdes dans les espaces, définition
\ref{spaces-groupoids-definition-quotient-sheaf})
est un espace algébrique, le morphisme $X \to X/G$ est
surjectif, plat et localement de présentation finie, et
$X$ est un torseur fppf sous $G$ au-dessus de $X/G$.
\end{lemma}

\begin{proof}
Le fait que $X/G$ soit un espace algébrique résulte immédiatement du
théorème \ref{theorem-final-bootstrap}
et des définitions. En effet, $X/G = X/R$, où $R = G \times_B X$.
Les morphismes $s, t : G \times_B X \to X$ sont plats et localement de
présentation finie (c'est clair pour $s$, qui est un changement de base de $G \to B$,
et cela en résulte pour $t$ par symétrie, en utilisant l'inverse), et le morphisme
$j : G \times_B X \to X \times_B X$ est un monomorphisme d'après
Groupoïdes dans les espaces, lemme \ref{spaces-groupoids-lemma-free-action},
puisque l'action est libre. Le morphisme $X \to X/G$ est surjectif, plat et
localement de présentation finie d'après le lemme \ref{lemma-covering-quotient}.
Pour voir que $X \to X/G$ est un torseur fppf sous $G$
(Groupoïdes dans les espaces, définition
\ref{spaces-groupoids-definition-principal-homogeneous-space})
il faut montrer que $G \times_S X \to X \times_{X/G} X$
est un isomorphisme et que $X \to X/G$ admet localement des sections pour la topologie fppf.
La seconde assertion résulte des propriétés déjà établies de $X \to X/G$.
Le morphisme $G \times_S X \to X \times_{X/G} X$ est injectif (comme morphisme de
faisceaux fppf), puisque l'action est libre. Enfin, il est aussi surjectif
comme morphisme de faisceaux d'après Groupoïdes dans les espaces, lemme
\ref{spaces-groupoids-lemma-quotient-pre-equivalence}.
Cela achève la démonstration.
\end{proof}

\begin{lemma}
\label{lemma-descent-torsor}
Soit $\{S_i \to S\}_{i \in I}$ un recouvrement de $(\Sch/S)_{fppf}$.
Soit $G$ un espace algébrique en groupes sur $S$, et notons
$G_i = G_{S_i}$ ses changements de base. Supposons donnés
\begin{enumerate}
\item pour chaque $i \in I$, un torseur fppf sous $G_i$, noté $X_i$, au-dessus de $S_i$ ;
\item pour chaque $i, j \in I$, un isomorphisme $G_{S_i \times_S S_j}$-équivariant
$\varphi_{ij} : X_i \times_S S_j \to S_i \times_S X_j$ satisfaisant à la condition
de cocycle sur chaque $S_i \times_S S_j \times_S S_j$.
\end{enumerate}
Alors il existe un torseur fppf sous $G$, noté $X$, au-dessus de $S$
dont le changement de base à $S_i$ est isomorphe à $X_i$, de telle sorte que l'on
retrouve la donnée de descente $\varphi_{ij}$.
\end{lemma}

\begin{proof}
Nous pouvons considérer $X_i$ comme un faisceau sur $(\Sch/S_i)_{fppf}$, voir
Espaces, section \ref{spaces-section-change-base-scheme}.
D'après
Sites, section \ref{sites-section-glueing-sheaves},
la donnée de descente $(X_i, \varphi_{ij})$ est effective au sens où
il existe un unique faisceau $X$ sur $(\Sch/S)_{fppf}$ qui
redonne les espaces algébriques $X_i$ après restriction à
$(\Sch/S_i)_{fppf}$. On a donc
$X_i = h_{S_i} \times X$. D'après le
lemme \ref{lemma-locally-algebraic-space},
$X$ est un espace algébrique, sous réserve de vérifier que $\coprod X_i$
est un espace algébrique, ce que nous ferons à la fin de la démonstration.
Par l'équivalence de catégories de
Sites, lemme \ref{sites-lemma-mapping-property-glue},
les morphismes d'action $G_i \times_{S_i} X_i \to X_i$
se recollent en un morphisme $a : G \times_S X \to X$.
Il faut maintenant montrer que $a$ est une action, que $X$
est un pseudo-torseur et qu'il est localement trivial pour la topologie fppf (voir
Groupoïdes dans les espaces,
définition \ref{spaces-groupoids-definition-principal-homogeneous-space}).
Ces propriétés se vérifient localement pour la topologie fppf et
découlent donc des propriétés correspondantes des actions
$G_i \times_{S_i} X_i \to X_i$. Le lemme est ainsi démontré.

\medskip\noindent
Nous suggérons au lecteur de passer la suite de la démonstration, qui est purement
ensembliste. Choisissons des recouvrements $\{S_{ij} \to S_j\}_{j \in J_i}$ de
$(\Sch/S)_{fppf}$
qui trivialisent les torseurs sous $G_i$ notés $X_i$ (c'est possible par hypothèse et par
Topologies, lemme \ref{topologies-lemma-fppf-induced}, partie (1)).
Alors $\{S_{ij} \to S\}_{i \in I, j \in J_i}$ est un recouvrement de
$(\Sch/S)_{fppf}$ ; nous pouvons donc supposer que chaque $X_i$
est le torseur trivial ! Bien entendu, nous pouvons encore raffiner le recouvrement ;
nous pouvons donc supposer que chaque $S_i$ est affine et que le cardinal de l'ensemble
d'indices $I$ est borné par le cardinal de l'ensemble des points
de $S$. Choisissons $U \in \Ob((\Sch/S)_{fppf})$ et un morphisme
étale surjectif $U \to G$. Alors $U_i = U \times_S S_i$ est muni
d'un morphisme étale surjectif vers $X_i \cong G_i$. D'après
Ensembles, lemme \ref{sets-lemma-bound-size-fibre-product},
on a $\text{size}(U_i) \leq \max\{\text{size}(U), \text{size}(S_i)\}$. D'après
Ensembles, lemme \ref{sets-lemma-bound-finite-type},
on a $\text{size}(S_i) \leq \text{size}(S)$.
Ainsi,
$\text{size}(U_i) \leq \max\{\text{size}(U), \text{size}(S)\}$
pour tout $i \in I$. En combinant cela avec la borne sur $|I|$ obtenue plus haut,
Ensembles, lemme \ref{sets-lemma-bound-size},
donne $\text{size}(\coprod U_i) \leq \max\{\text{size}(U), \text{size}(S)\}$.
Par conséquent,
Espaces, lemme \ref{spaces-lemma-coproduct-algebraic-spaces},
s'applique et montre que $\coprod X_i$ est un espace algébrique, comme
il fallait le prouver.
\end{proof}




\section{Espaces algébriques pour la topologie étale}
\label{section-spaces-etale}

\noindent
Soit $S$ un schéma. Au lieu de travailler avec des faisceaux sur
le grand site fppf $(\Sch/S)_{fppf}$, nous pourrions travailler avec des faisceaux
sur le grand site étale $(\Sch/S)_\etale$. Tous les résultats de
Espaces algébriques, sections \ref{spaces-section-representable} et
\ref{spaces-section-representable-properties}
ont un sens pour les faisceaux sur $(\Sch/S)_\etale$.
On obtient ainsi une seconde notion d'espace algébrique en travaillant pour la
topologie étale. Cette notion est a priori plus faible que celle introduite dans
Espaces algébriques, définition \ref{spaces-definition-algebraic-space},
puisqu'un faisceau pour la topologie fppf est assurément un faisceau pour la
topologie étale. Les deux notions sont toutefois équivalentes, comme le montre le
lemme suivant.

\begin{lemma}
\label{lemma-spaces-etale}
La catégorie sous-jacente commune à $\Sch_{fppf}$ et $\Sch_\etale$ sera notée
$\Sch_\alpha$
(voir Topologies, remarque \ref{topologies-remark-choice-sites}).
Soit $S$ un objet de $\Sch_\alpha$. Soit
$$
F : (\Sch_\alpha/S)^{opp} \longrightarrow \textit{Sets}
$$
un préfaisceau possédant les propriétés suivantes :
\begin{enumerate}
\item $F$ est un faisceau pour la topologie étale ;
\item la diagonale $\Delta : F \to F \times F$ est représentable ;
\item il existe $U \in \Ob(\Sch_\alpha/S)$
et un morphisme $U \to F$ surjectif et étale.
\end{enumerate}
Alors $F$ est un espace algébrique au sens de
Espaces algébriques, définition \ref{spaces-definition-algebraic-space}.
\end{lemma}

\begin{proof}
Notons que les propriétés (2) et (3) du lemme et les propriétés correspondantes
(2) et (3) de
Espaces algébriques, définition \ref{spaces-definition-algebraic-space},
ne dépendent pas de la topologie. En effet, ces propriétés
ne font intervenir que la notion de produit fibré de préfaisceaux, les morphismes de
préfaisceaux, la notion de transformation représentable de foncteurs
et ce que signifie pour une telle transformation d'être surjective et étale.
Il ne reste donc qu'à prouver qu'un faisceau étale $F$ possédant les propriétés
(2) et (3) est aussi un faisceau fppf.

\medskip\noindent
Pour cela, posons $R = U \times_F U$. D'après (2), le préfaisceau $R$ est représenté
par un schéma et, d'après (3), les projections $R \to U$ sont étales. Ainsi,
$j : R \to U \times_S U$ est une relation d'équivalence étale. De plus,
$U \to F$ identifie $F$ au quotient de $U$ par $R$ pour la
topologie étale : (a) si $T \to F$ est un morphisme, alors $\{T \times_F U \to T\}$
est un recouvrement étale, donc $U \to F$ est un épimorphisme de faisceaux pour la
topologie étale ; (b) si $a, b : T \to U$ ont pour image la même section de $F$,
alors $(a, b) : T \to R$, donc $a$ et $b$ ont la même image dans le quotient
de $U$ par $R$ pour la topologie étale. Notons ensuite $U/R$ le faisceau quotient
pour la topologie fppf ; c'est un espace algébrique d'après
Espaces, théorème \ref{spaces-theorem-presentation}.
Nous avons donc des morphismes (transformations de foncteurs)
$$
U \to F \to U/R.
$$
D'après le
théorème \ref{spaces-theorem-presentation} d'Espaces cité plus haut,
le composé est représentable, surjectif et étale. Ainsi, pour tout
schéma $T$ et tout morphisme $T \to U/R$, le produit fibré $V = T \times_{U/R} U$
est un schéma étale et surjectif sur $T$. Autrement dit, $\{V \to U\}$
est un recouvrement étale. Cela prouve que $U \to U/R$ est surjectif comme
morphisme de faisceaux pour la topologie étale. Il s'ensuit que
$F \to U/R$ est surjectif comme morphisme de faisceaux pour la topologie étale.
D'autre part, le morphisme $F \to U/R$ est injectif (comme morphisme de préfaisceaux),
puisque $R = U \times_{U/R} U$, de nouveau d'après
Espaces, théorème \ref{spaces-theorem-presentation}.
Ainsi, $F \to U/R$ est un isomorphisme de faisceaux étales, voir
Sites, lemme \ref{sites-lemma-mono-epi-sheaves},
ce qui achève la démonstration.
\end{proof}

\noindent
Il existe aussi un analogue de
Espaces, lemme \ref{spaces-lemma-etale-locally-representable-gives-space}.

\begin{lemma}
\label{lemma-spaces-etale-locally-representable}
La catégorie sous-jacente commune à $\Sch_{fppf}$ et $\Sch_\etale$ sera notée
$\Sch_\alpha$
(voir Topologies, remarque \ref{topologies-remark-choice-sites}).
Soit $S$ un objet de $\Sch_\alpha$. Soit
$$
F : (\Sch_\alpha/S)^{opp} \longrightarrow \textit{Sets}
$$
un préfaisceau possédant les propriétés suivantes :
\begin{enumerate}
\item $F$ est un faisceau pour la topologie étale ;
\item il existe un espace algébrique $U$ sur $S$
et un morphisme $U \to F$ représentable par
des espaces algébriques, surjectif et étale.
\end{enumerate}
Alors $F$ est un espace algébrique au sens de
Espaces algébriques, définition \ref{spaces-definition-algebraic-space}.
\end{lemma}

\begin{proof}
Posons $R = U \times_F U$. C'est un espace algébrique puisque $U \to F$ est supposé
représentable par des espaces algébriques. Les projections $s, t : R \to U$ sont des
morphismes étales d'espaces algébriques puisque $U \to F$ est supposé étale.
Le morphisme $j = (t, s) : R \to U \times_S U$ est un monomorphisme et une
relation d'équivalence puisque $R = U \times_F U$. D'après le
théorème \ref{theorem-final-bootstrap},
le faisceau quotient fppf $F' = U/R$ est un espace algébrique.
Le morphisme $U \to F'$ est surjectif, plat et localement de présentation
finie d'après le lemme \ref{lemma-covering-quotient}.
Le morphisme $R \to U \times_{F'} U$ est surjectif comme morphisme de faisceaux
fppf d'après Groupoïdes dans les espaces, lemme
\ref{spaces-groupoids-lemma-quotient-pre-equivalence}
et, puisque $j$ est un monomorphisme, c'est un isomorphisme.
Ainsi, le changement de base de $U \to F'$ par $U \to F'$ est étale,
et l'on en déduit que $U \to F'$ est étale d'après
Descente sur les espaces, lemme \ref{spaces-descent-lemma-descending-property-etale}.
Le morphisme $U \to F'$ est donc surjectif comme morphisme de faisceaux étales.
Cela signifie que $F'$ est égal au faisceau quotient $U/R$
pour la topologie étale (petite vérification omise). On obtient donc
une factorisation canonique $U \to F' \to F$, et $F' \to F$ est un morphisme
injectif de faisceaux. D'autre part, $U \to F$ est surjectif comme morphisme
de faisceaux étales, et il en va donc de même de $F' \to F$. Ainsi $F' = F$,
ce qui achève la démonstration.
\end{proof}

\noindent
En fait, il suffit de disposer d'un recouvrement lisse par un schéma et de supposer
que la diagonale est représentable par des espaces algébriques.

\begin{lemma}
\label{lemma-spaces-etale-smooth-cover}
La catégorie sous-jacente commune à $\Sch_{fppf}$
et $\Sch_\etale$ sera notée $\Sch_\alpha$ (voir
Topologies, remarque \ref{topologies-remark-choice-sites}). Soit $S$ un objet
de $\Sch_\alpha$. Soit
$$
F : (\Sch_\alpha/S)^{opp} \longrightarrow \textit{Sets}
$$
un préfaisceau possédant les propriétés suivantes :
\begin{enumerate}
\item $F$ est un faisceau pour la topologie étale ;
\item la diagonale $\Delta : F \to F \times F$ est représentable
par des espaces algébriques ;
\item il existe $U \in \Ob(\Sch_\alpha/S)$
et un morphisme $U \to F$ surjectif et lisse.
\end{enumerate}
Alors $F$ est un espace algébrique au sens de
Espaces algébriques, définition \ref{spaces-definition-algebraic-space}.
\end{lemma}

\begin{proof}
La démonstration suit celle du lemme \ref{lemma-spaces-etale}. Posons
$R = U \times_F U$. D'après (2), le préfaisceau $R$ est un espace algébrique et, d'après (3),
les projections $R \to U$ sont lisses et surjectives. Notons $(U, R, s, t, c)$
le groupoïde associé à la relation d'équivalence $j : R \to U \times_S U$
(voir Groupoïdes dans les espaces, lemme
\ref{spaces-groupoids-lemma-equivalence-groupoid}).
D'après le théorème \ref{theorem-final-bootstrap}, $X = U/R$ (quotient
pour la topologie fppf) est un espace algébrique. Comme la topologie lisse
et la topologie étale ont les mêmes faisceaux (d'après
Compléments sur les morphismes, lemme \ref{more-morphisms-lemma-etale-dominates-smooth}),
le morphisme $U \to F$ identifie $F$ au quotient de
$U$ par $R$ pour la topologie lisse (détails omis).
Nous avons donc des morphismes (transformations de foncteurs)
$$
U \to F \to X.
$$
D'après le lemme \ref{lemma-covering-quotient}, $U \to X$ est
surjectif, plat et localement de présentation finie. D'après
Groupoïdes dans les espaces, lemme
\ref{spaces-groupoids-lemma-quotient-pre-equivalence}
(et puisque $j$ est un monomorphisme), on a $R = U \times_X U$. D'après
Descente sur les espaces, lemme \ref{spaces-descent-lemma-descending-property-smooth},
on en déduit que $U \to X$ est lisse et surjectif (puisque les projections
$R \to U$ sont lisses et surjectives et que $\{U \to X\}$ est un recouvrement
fppf). Ainsi, pour tout schéma $T$ et tout morphisme $T \to X$, le produit fibré
$T \times_X U$ est un espace algébrique lisse et surjectif sur $T$.
Choisissons un schéma $V$ et un morphisme étale surjectif $V \to T \times_X U$.
Alors $\{V \to T\}$ est un recouvrement lisse tel que $V \to T \to X$
se relève en un morphisme $V \to U$. Cela prouve que
$U \to X$ est surjectif comme morphisme de faisceaux pour la topologie lisse.
Il s'ensuit que $F \to X$ est surjectif comme morphisme de faisceaux pour la topologie
lisse. D'autre part, le morphisme $F \to X$ est injectif (comme morphisme
de préfaisceaux), puisque $R = U \times_X U$.
Ainsi, $F \to X$ est un isomorphisme de faisceaux lisses ($=$ étales),
voir Sites, lemme \ref{sites-lemma-mono-epi-sheaves},
ce qui achève la démonstration.
\end{proof}

\noindent
Enfin, voici l'analogue de
Espaces, lemme \ref{spaces-lemma-etale-locally-representable-gives-space},
où l'espace est recouvert par un morphisme lisse.

\begin{lemma}
\label{lemma-spaces-smooth-locally-representable}
La catégorie sous-jacente commune à $\Sch_{fppf}$ et $\Sch_\etale$ sera notée
$\Sch_\alpha$
(voir Topologies, remarque \ref{topologies-remark-choice-sites}).
Soit $S$ un objet de $\Sch_\alpha$. Soit
$$
F : (\Sch_\alpha/S)^{opp} \longrightarrow \textit{Sets}
$$
un préfaisceau possédant les propriétés suivantes :
\begin{enumerate}
\item $F$ est un faisceau pour la topologie étale ;
\item il existe un espace algébrique $U$ sur $S$
et un morphisme $U \to F$ représentable par
des espaces algébriques, surjectif et lisse.
\end{enumerate}
Alors $F$ est un espace algébrique au sens de
Espaces algébriques, définition \ref{spaces-definition-algebraic-space}.
\end{lemma}

\begin{proof}
La démonstration est identique à celle du
lemme \ref{lemma-spaces-etale-locally-representable}.
Posons $R = U \times_F U$. C'est un espace algébrique puisque $U \to F$ est supposé
représentable par des espaces algébriques. Les projections $s, t : R \to U$ sont des
morphismes lisses d'espaces algébriques puisque $U \to F$ est supposé lisse.
Le morphisme $j = (t, s) : R \to U \times_S U$ est un monomorphisme et une
relation d'équivalence puisque $R = U \times_F U$. D'après le
théorème \ref{theorem-final-bootstrap},
le faisceau quotient fppf $F' = U/R$ est un espace algébrique.
Le morphisme $U \to F'$ est surjectif, plat et localement de présentation
finie d'après le lemme \ref{lemma-covering-quotient}.
Le morphisme $R \to U \times_{F'} U$ est surjectif comme morphisme de faisceaux
fppf d'après Groupoïdes dans les espaces, lemme
\ref{spaces-groupoids-lemma-quotient-pre-equivalence}
et, puisque $j$ est un monomorphisme, c'est un isomorphisme.
Ainsi, le changement de base de $U \to F'$ par $U \to F'$ est lisse,
et l'on en déduit que $U \to F'$ est lisse d'après
Descente sur les espaces, lemme \ref{spaces-descent-lemma-descending-property-smooth}.
Le morphisme $U \to F'$ est donc surjectif comme morphisme de faisceaux étales (car
la topologie lisse et la topologie étale ont les mêmes faisceaux d'après
Compléments sur les morphismes, lemme \ref{more-morphisms-lemma-etale-dominates-smooth}).
Cela signifie que $F'$ est égal au faisceau quotient $U/R$
pour la topologie étale (petite vérification omise). On obtient donc
une factorisation canonique $U \to F' \to F$, et $F' \to F$ est un morphisme
injectif de faisceaux. D'autre part, $U \to F$ est surjectif comme morphisme
de faisceaux étales (puisque les topologies lisse et étale ont les mêmes
faisceaux), et il en va donc de même de $F' \to F$. Ainsi $F' = F$,
ce qui achève la démonstration.
\end{proof}






\begin{multicols}{2}[\section{Autres chapitres}]
\noindent
Préliminaires
\begin{enumerate}
\item \hyperref[introduction-section-phantom]{Introduction}
\item \hyperref[conventions-section-phantom]{Conventions}
\item \hyperref[sets-section-phantom]{Théorie des ensembles}
\item \hyperref[categories-section-phantom]{Catégories}
\item \hyperref[topology-section-phantom]{Topologie}
\item \hyperref[sheaves-section-phantom]{Faisceaux sur les espaces}
\item \hyperref[sites-section-phantom]{Sites et faisceaux}
\item \hyperref[stacks-section-phantom]{Champs}
\item \hyperref[fields-section-phantom]{Corps}
\item \hyperref[algebra-section-phantom]{Algèbre commutative}
\item \hyperref[brauer-section-phantom]{Groupes de Brauer}
\item \hyperref[homology-section-phantom]{Algèbre homologique}
\item \hyperref[derived-section-phantom]{Catégories dérivées}
\item \hyperref[simplicial-section-phantom]{Méthodes simpliciales}
\item \hyperref[more-algebra-section-phantom]{Compléments d'algèbre}
\item \hyperref[smoothing-section-phantom]{Lissage des morphismes d'anneaux}
\item \hyperref[modules-section-phantom]{Faisceaux de modules}
\item \hyperref[sites-modules-section-phantom]{Modules sur les sites}
\item \hyperref[injectives-section-phantom]{Objets injectifs}
\item \hyperref[cohomology-section-phantom]{Cohomologie des faisceaux}
\item \hyperref[sites-cohomology-section-phantom]{Cohomologie sur les sites}
\item \hyperref[dga-section-phantom]{Algèbre différentielle graduée}
\item \hyperref[dpa-section-phantom]{Algèbre à puissances divisées}
\item \hyperref[sdga-section-phantom]{Faisceaux différentiels gradués}
\item \hyperref[hypercovering-section-phantom]{Hyperrecouvrements}
\end{enumerate}
Schémas
\begin{enumerate}
\setcounter{enumi}{25}
\item \hyperref[schemes-section-phantom]{Schémas}
\item \hyperref[constructions-section-phantom]{Constructions de schémas}
\item \hyperref[properties-section-phantom]{Propriétés des schémas}
\item \hyperref[morphisms-section-phantom]{Morphismes de schémas}
\item \hyperref[coherent-section-phantom]{Cohomologie des schémas}
\item \hyperref[divisors-section-phantom]{Diviseurs}
\item \hyperref[limits-section-phantom]{Limites de schémas}
\item \hyperref[varieties-section-phantom]{Variétés}
\item \hyperref[topologies-section-phantom]{Topologies sur les schémas}
\item \hyperref[descent-section-phantom]{Descente}
\item \hyperref[perfect-section-phantom]{Catégories dérivées de schémas}
\item \hyperref[more-morphisms-section-phantom]{Compléments sur les morphismes}
\item \hyperref[flat-section-phantom]{Compléments sur la platitude}
\item \hyperref[groupoids-section-phantom]{Schémas en groupoïdes}
\item \hyperref[more-groupoids-section-phantom]{Compléments sur les schémas en groupoïdes}
\item \hyperref[etale-section-phantom]{Morphismes étales de schémas}
\end{enumerate}
Thèmes de théorie des schémas
\begin{enumerate}
\setcounter{enumi}{41}
\item \hyperref[chow-section-phantom]{Homologie de Chow}
\item \hyperref[intersection-section-phantom]{Théorie de l'intersection}
\item \hyperref[pic-section-phantom]{Schémas de Picard des courbes}
\item \hyperref[weil-section-phantom]{Théories cohomologiques de Weil}
\item \hyperref[adequate-section-phantom]{Modules adéquats}
\item \hyperref[dualizing-section-phantom]{Complexes dualisants}
\item \hyperref[duality-section-phantom]{Dualité pour les schémas}
\item \hyperref[discriminant-section-phantom]{Discriminants et différentes}
\item \hyperref[derham-section-phantom]{Cohomologie de de Rham}
\item \hyperref[local-cohomology-section-phantom]{Cohomologie locale}
\item \hyperref[algebraization-section-phantom]{Géométrie algébrique et formelle}
\item \hyperref[curves-section-phantom]{Courbes algébriques}
\item \hyperref[resolve-section-phantom]{Résolution des surfaces}
\item \hyperref[models-section-phantom]{Réduction semi-stable}
\item \hyperref[functors-section-phantom]{Foncteurs et morphismes}
\item \hyperref[equiv-section-phantom]{Catégories dérivées de variétés}
\item \hyperref[pione-section-phantom]{Groupes fondamentaux de schémas}
\item \hyperref[etale-cohomology-section-phantom]{Cohomologie étale}
\item \hyperref[crystalline-section-phantom]{Cohomologie cristalline}
\item \hyperref[proetale-section-phantom]{Cohomologie pro-étale}
\item \hyperref[relative-cycles-section-phantom]{Cycles relatifs}
\item \hyperref[more-etale-section-phantom]{Compléments de cohomologie étale}
\item \hyperref[trace-section-phantom]{La formule des traces}
\end{enumerate}
Espaces algébriques
\begin{enumerate}
\setcounter{enumi}{64}
\item \hyperref[spaces-section-phantom]{Espaces algébriques}
\item \hyperref[spaces-properties-section-phantom]{Propriétés des espaces algébriques}
\item \hyperref[spaces-morphisms-section-phantom]{Morphismes d'espaces algébriques}
\item \hyperref[decent-spaces-section-phantom]{Espaces algébriques décents}
\item \hyperref[spaces-cohomology-section-phantom]{Cohomologie des espaces algébriques}
\item \hyperref[spaces-limits-section-phantom]{Limites d'espaces algébriques}
\item \hyperref[spaces-divisors-section-phantom]{Diviseurs sur les espaces algébriques}
\item \hyperref[spaces-over-fields-section-phantom]{Espaces algébriques sur des corps}
\item \hyperref[spaces-topologies-section-phantom]{Topologies sur les espaces algébriques}
\item \hyperref[spaces-descent-section-phantom]{Descente et espaces algébriques}
\item \hyperref[spaces-perfect-section-phantom]{Catégories dérivées d'espaces}
\item \hyperref[spaces-more-morphisms-section-phantom]{Compléments sur les morphismes d'espaces}
\item \hyperref[spaces-flat-section-phantom]{Platitude sur les espaces algébriques}
\item \hyperref[spaces-groupoids-section-phantom]{Groupoïdes dans les espaces algébriques}
\item \hyperref[spaces-more-groupoids-section-phantom]{Compléments sur les groupoïdes dans les espaces}
\item \hyperref[bootstrap-section-phantom]{Amorçage}
\item \hyperref[spaces-pushouts-section-phantom]{Sommes amalgamées d'espaces algébriques}
\end{enumerate}
Thèmes de géométrie
\begin{enumerate}
\setcounter{enumi}{81}
\item \hyperref[spaces-chow-section-phantom]{Groupes de Chow des espaces}
\item \hyperref[groupoids-quotients-section-phantom]{Quotients de groupoïdes}
\item \hyperref[spaces-more-cohomology-section-phantom]{Compléments sur la cohomologie des espaces}
\item \hyperref[spaces-simplicial-section-phantom]{Espaces simpliciaux}
\item \hyperref[spaces-duality-section-phantom]{Dualité pour les espaces}
\item \hyperref[formal-spaces-section-phantom]{Espaces algébriques formels}
\item \hyperref[restricted-section-phantom]{Algébrisation des espaces formels}
\item \hyperref[spaces-resolve-section-phantom]{Résolution des surfaces revisitée}
\end{enumerate}
Théorie des déformations
\begin{enumerate}
\setcounter{enumi}{89}
\item \hyperref[formal-defos-section-phantom]{Théorie formelle des déformations}
\item \hyperref[defos-section-phantom]{Théorie des déformations}
\item \hyperref[cotangent-section-phantom]{Le complexe cotangent}
\item \hyperref[examples-defos-section-phantom]{Problèmes de déformation}
\end{enumerate}
Champs algébriques
\begin{enumerate}
\setcounter{enumi}{93}
\item \hyperref[algebraic-section-phantom]{Champs algébriques}
\item \hyperref[examples-stacks-section-phantom]{Exemples de champs}
\item \hyperref[stacks-sheaves-section-phantom]{Faisceaux sur les champs algébriques}
\item \hyperref[criteria-section-phantom]{Critères de représentabilité}
\item \hyperref[artin-section-phantom]{Axiomes d'Artin}
\item \hyperref[quot-section-phantom]{Espaces de Quot et de Hilbert}
\item \hyperref[stacks-properties-section-phantom]{Propriétés des champs algébriques}
\item \hyperref[stacks-morphisms-section-phantom]{Morphismes de champs algébriques}
\item \hyperref[stacks-limits-section-phantom]{Limites de champs algébriques}
\item \hyperref[stacks-cohomology-section-phantom]{Cohomologie des champs algébriques}
\item \hyperref[stacks-perfect-section-phantom]{Catégories dérivées de champs}
\item \hyperref[stacks-introduction-section-phantom]{Introduction aux champs algébriques}
\item \hyperref[stacks-more-morphisms-section-phantom]{Compléments sur les morphismes de champs}
\item \hyperref[stacks-geometry-section-phantom]{La géométrie des champs}
\end{enumerate}
Thèmes de théorie des espaces de modules
\begin{enumerate}
\setcounter{enumi}{107}
\item \hyperref[moduli-section-phantom]{Champs de modules}
\item \hyperref[moduli-curves-section-phantom]{Espaces de modules de courbes}
\end{enumerate}
Divers
\begin{enumerate}
\setcounter{enumi}{109}
\item \hyperref[examples-section-phantom]{Exemples}
\item \hyperref[exercises-section-phantom]{Exercices}
\item \hyperref[guide-section-phantom]{Guide bibliographique}
\item \hyperref[desirables-section-phantom]{Desiderata}
\item \hyperref[coding-section-phantom]{Style de programmation}
\item \hyperref[obsolete-section-phantom]{Obsolète}
\item \hyperref[fdl-section-phantom]{Licence de documentation libre GNU}
\item \hyperref[index-section-phantom]{Index généré automatiquement}
\end{enumerate}
\end{multicols}



\bibliography{my}
\bibliographystyle{amsalpha}

\end{document}
